\quellseite{587}

\part{Das erste Buch der Chronika}

\chapter{Das 1. Kapitel}

\section[Geschlechtsregister der Patriarchen bis auf Jakob]{Geschlechtsregister der Patriarchen bis auf Jakob}
\quellseite{588}

\chapter{Das 2. Capitel}

\section[Söhne Jakobs und Juda]{Söhne Jakobs und Juda}
\quellseite{589}

a Seher, die Propheten heißen Seher oder Schauer in der Schrift, und sehen heißt lehren oder predigen. Also ist Sobal ein Prophet oder Prediger gewesen in dem halben Manuhoth, das ist ein Name eines Ländleins, das vielleicht fein und wohl gelegen ist; darum es Manuhoth, das ist, Ruhe heißt.

\chapter{Das 3. Capitel}

\section[Register der Könige Juda]{Register der Könige Juda}

a Egla diese wird auch Davids Weib genennet, vielleicht, daß sie die erste und einige ist gewesen, da er noch der Schafe gehütet hat.
\quellseite{590}

\chapter{Das 4. Kapitel}

\section[Nachkommen Juda]{Nachkommen Juda}
\quellseite{591}

\chapter{Das 5. Kapitel}

\section[Geschlechter Simeons]{Geschlechter Simeons}

\chapter{Das 6. (sonst 5.) Kapitel}

\section[Von Ruben, Gad, Manasse]{Von Ruben, Gad, Manasse}
\quellseite{592}

\chapter{Das 7. (Sonst 6.) Capitel}

\section[Zahl und Wohnung der Kinder Levi und Aaron]{Zahl und Wohnung der Kinder Levi und Aaron}

a Erbitten siehe, wie der Glaube alles vermag, wie Christus spricht Matth. 21, 22.

b Nämlich durch Salmanasser.
% Scan page 593: omitted; Bible text/narrative only, no retained apparatus.
\quellseite{594}

\chapter{Das 8. (sonst 7.) Kapitel}

\section[Geschlechtsregister der sechs übrigen Stämme Jakobs]{Geschlechtsregister der sechs übrigen Stämme Jakobs}
% Scan page 595: omitted; Bible text/narrative only, no retained apparatus.
\quellseite{596}

\chapter{Das 9. (sonst 8.) Capitel}

\section[Das Haus Benjamin und Saul]{Das Haus Benjamin und Saul}

\chapter{Das 10. (sonst 9.) Capitel}

\section[Einwohner zu Jerusalem und Gibeon]{Einwohner zu Jerusalem und Gibeon}

a Nethinim heißt die gegebenen oder geschenkten, und waren die Leviten, die sich insonderheit zum Gottesdienst begeben hatten, wo Gottes Hütten und die Lade war unter den Priestern.
\quellseite{597}

b Durch ihren Glauben) das ist gesagt: Solche Stifte thaten sie nicht aus menschlicher Andacht und guter Meinung, sondern aus Gottes Befehl und ihrem Glauben. Denn in Gottes Sache soll keine menschliche Stiftung oder Werk gelten.
\quellseite{598}

\chapter{Das 11. (sonst 10.) Kapitel}

\section{Vorrede}

Von Strafen derer, so Gottes Wort verachten und dem ungehorsam sein, sagt Gott im Propheten Hosea am 4. Kapitel V. 6.: Mein Volk ist dahin, darum, daß es nicht lernen will. Denn du verwirfst Gottes Wort, darum will ich dich auch verwerfen. Du vergissest des Gesetzes deines Gottes, darum will ich auch deiner Kinder vergessen. Und in Sprüchen Salomo im 1. Kapitel V. 29. sagt Gott abermals: Darum, daß sie hassten die Lehre, und wollten des Herrn Furcht nicht haben, wollten meines Raths nicht, und lästerten alle meine Strafe; so sollen sie essen von den Früchten ihres Wesens, und ihres Raths satt werden. Dagegen von Belohnung derer, die Gottes Wort hoch achten und dem folgen, sagt Gott Esaia 55. Kapitel V. 2. 3. also: Höret mir doch zu, und esset das Gute, so wird eure Seele in Wollust fett werden. Neiget eure Ohren her, und kommt her zu mir; höret, so wird eure Seele leben: denn ich will mit euch einen ewigen Bund machen. Und abermals in Sprüchen Salomo 3. Kapitel V. 21. 22. spricht er: Laß sie (die Weisheit) nicht von deinen Augen weichen, daß wird deiner Seele Leben sein. Dies beides lehret auch das 11. Kapitel des ersten Buchs der Chronika.

\section[Sauls Untergang im Streit wider die Philister]{Sauls Untergang im Streit wider die Philister}
\quellseite{599}

\section{Summarium}

Dies ist eine Historie, wie es mit dem König Saul habe ein Ende genommen, darum, daß er erstlich an Gottes Wort nicht gehalten, sondern wider Gottes Befehl mit den Amalekitern gehandelt und hernach die Wahrsagerin, und nicht des HErrn Rath gefraget hatte, V. 5. Solche beide Sünden scheinen, als wären es geringe Sünden, aber die Strafe zeiget genugsam an, daß es große, schwere Sünden sind, sintemal Gott dieselbigen an allen seinen Nachkommen gestrafet hat. Darum soll jedermann allein auf Gottes Wort sehen, und demselbigen mit Fleiß folgen, und nichts dawider fürnehmen, sonderlich aber die Obrigkeit. Sonst soll es ihnen gehen, wie es Saul und seinem Geschlecht gegangen.

\section{Votum}

GOtt regiere und führe unsere Obrigkeit und uns alle, daß wir seinem Wort folgen und wider seinen Befehl nichts fürnehmen noch handeln; und lasse ihm uns und die Unsern, auch insonderheit alle christliche Regenten, in Gnaden befohlen sein, Amen.

\chapter{Das 12. (sonst 11.) Kapitel}

\section{Vorrede}

Gleich wie die Ochsen nicht von Ochsen, und die Schafe nicht von Schafen regieret werden, sondern von einer edlen Kreatur, nämlich von Menschen: also ist die menschliche Natur schwächer, denn daß sie sich selber regieren kann, das zu lassen und zu thun, was Gott geboten und verboten hat, zu seiner Ehre und zu Nutz den Nebenmenschen. Derwegen hat Gott Land und Leute zu Regenten fürgestellet, nicht gemeine Leute, sondern göttliche Menschen, nämlich die Heroas und Wunderleute, derer Regierung (denn dazu gibt Gott ihnen Glück) gottliebenden und gehorsamen Unterthanen heilsam und nützlich ist. Daher werden sie genennet Väter des Vaterlandes, Schilde der Erden und Bäume, von denen und unter denen, die Menschen Schatten, Schutz und Hülfe haben. Dieser Lehre Exempel hält uns für das 12. Kapitel.

\section[Salbung Davids zum Könige. Seine Wohnung und streitbaren]{Salbung Davids zum Könige. Seine Wohnung und streitbaren Helden}

a (Dreihundert) 2 Sam. 23, 8. stehet achthundert. Wer einen Haufen von achthundert angreift, und schlägt dreihundert todt, und die fünfhundert in die Flucht, der hat sie alle auf die Flucht geschlagen. Also da hie (unter dreißigen) steht, hat droben 2 Sam. 23, 8. der Text unter dreien. Denn welcher unter den besten dreien der Fürnehmste ist, der ist auch unter den dreißigen allen der Fürnehmste.
\quellseite{600}

b Goß das ist, er opferte es Gott zum Trankopfer.

\section{Summarium}

Ist eine Historia, wie David zum Könige erwählt wird, und ihm Gott die Gnade thut, daß er viele tapfere Helden an seinem Hofe hat, v. 3. 10. Denn wo Gott zum Regiment will Glück geben, da gibt er seine Leute, die zu Sachen geschickt und tüchtig sind. Solche Leute sollen die Regenten lieb und werth haben. Darum reuet
\quellseite{601}
es David, daß aus seinem Ausschicken die drey tapferen Helden Jasabeam, Eleasar und Samma ihr Leben so gering gewagt und in so große Gefahr gesetzt hatten, daß sie durch der Philister Lager gelaufen und David das Wasser aus dem Brunnen zu Bethlehem gebracht hatten. Und ist wahr, muthwilliges Wagen gefällt Gott nicht. Gleich wie Gott sonderlich Glück dazu gibt, wo man in der Noth auf seine Hülfe etwas waget. Aber wo man sich in Gefahr begibt, der man wohl geraten kann, da ist selten Glück bei.

\section{Votum}

Unser lieber GOtt und HErr lasse ihm unsere Obrigkeit auch befohlen sein, gebe und erhalte ihnen treue Räthe und Mitgehülffen; und erhalte bei uns sein heilig Wort und guten Frieden, und helfe, daß Obrigkeit und Unterthanen ihn allhie zeitlich und dort ewig loben und preisen, Amen.

\chapter{Das 13. (sonst 12.) Capitel}

\section{Davids Beistand an streitbaren Männern}

a Das ist, verborgen.
\quellseite{602}

b Kriegsverständig.

\chapter{Das 14. (sonst 13.) Kapitel}

\section{Vorrede}

Wiewohl der Obrigkeit Amt ist, aus Gottes Ordnung herrschen und regieren, Luc. 22, 25., so soll sie doch auch dienen; aber nicht dem Satan, nicht der Sünden, nicht Geld und Gut, nicht irdischer Wollust, sondern Gott und ihren Unterthanen soll sie dienen. Für ihre Person sollen alle die im Stande der Obrigkeit sein, Gott dienen, der sie erschaffen, durch seinen lieben Sohn erlöset, mit dem heil. Geist auf den Tag ihrer Erlösung versiegelt hat, Ephes. 4.; dem sie in der heiligen Taufe Gehorsam gelobet, und sollen ihm dienen, nach ihrer aufgerichteten Bestallung und Tabulatur der heiligen zehn Gebote Gottes. Darnach sollen sie auch ihren Unterthanen dienen in ihrem tragenden Amt, sollen dem großmächtigen Könige der Ehren ihre Thore, Land, Leute und Städte aufthun, Ps. 24, 7., sollen ihn hausen und beherbergen, Kirchen und Schulen treulich bestellen, mit höchstem Fleiß abschaffen alle Aergernisse in Lehre und Leben, nach Friede trachten, die Frommen schützen und die Bösen strafen, Ps. 34, 22. Ein solcher Regent ist gewesen Theodosius, der Kaiser, dem dieß Zeugniß gegeben wird, daß er sich in seiner Todesstunde ernstlicher und heftiger der christlichen Kirchen, als seines eigenen Leibes und Lebens angenommen hat. Ein solcher Regent ist auch gewesen der König David, von dem das 14. Kapitel.

\section[Ab- und Einholung der Lade des Bundes]{Ab- und Einholung der Lade des Bundes}
\quellseite{603}

\section{Summarium}

Ist ein fein fürstlich Exempel, wie der fromme König David sich mit Ernst darum annimmt, daß der Gottesdienst, so unter Sauls Regiment gefallen, wieder recht werde angerichtet, 2. 3. Denn die Obrigkeit soll den rechten Gottesdienst sonderlich fördern, wo sie anders wohl regieren und Gottes Segen haben will.

Daß aber der HErr den Usa tödtet, daß er die Lade anrühret (der doch von der Leviten Stamm nicht war), ist darum geschehen, daß Gott den ganzen Haufen mit solcher Strafe zur Furcht und Zucht hat gewöhnen wollen. Denn mit dem Gottesdienst soll niemand leichtfertig umgehen. Widerrum segnet Gott Obed-Edom, darum, daß er die Lade des HErrn beherberget.

\section{Votum}

Helffe GOtt, daß unsere liebe Obrigkeit die Gottesdienste auch ferner fördere, und verleihe, daß wir uns bei denen als wahre Christen erzeigen und halten; und segne reichlich, die in diesen betrübten Zeiten seiner Kirche und deren Gliedmaßen Herberge geben, tröstlich und beförderlich sein, Amen.

\chapter{Das 15. (sonst 14.) Kapitel}

\section{Vorrede}

Alle Dinge, die zu Gottes Ehren, und unserm Nächsten zu Nutz, getragen sollen, Anfang, Mittel und Ende, stehet in Gottes Händen, Gewalt und Segen. Derwegen wer regieren, predigen, handeln, obern, handthieren und arbeiten will, daß ihm gelinge und daß er Glück und Fortgang, dazu Nutz und Frommen davon habe, der fange es mit Gott an, mache auch mit Gott Feierabend, das ist, er bete, daß Gott den Anfang segne, dem Mittel beistehe und zum Ende befördere; also wird er erfahren, daß Gott ist ein Vater des Lichts, von dem alle gute Gaben von oben herab kommen, Jac. 1, 17. Und daß es umsonst ist, daß man frühe aufstehet, und darnach lange sitzet, und daß man das Brod isset mit Sorgen, daß Gott seinen Freunden solches schlafend gibt, Ps. 127, 3., das ist, daß er seinen Gläubigen und Gehorsamen über Nacht bescheret, und daß Gott mehr hab, denn er vergab; und da Tag kömmt, daß auch Rath kömmt, wie die alten Deutschen gesagt haben. Darauf höret auch an das 15. Kapitel.

\section[Davids Weiber, Kinder und Sieg wider die Philister]{Davids Weiber, Kinder und Sieg wider die Philister}
\quellseite{604}
(Baal-Prazim) Baal heißt ein Hauswirth oder Mann, Perez heißt ein Riß oder Fach, darum muß diese Stätte Baal Prazim heißen, weil die Philister da zertrennt und zerrissen sind.

\section{Summarium}

Dies ist eine Historie, wie David in allem zuvor GOtt rathfragend und sich seinem Wort nach hält, V. 10. Darum gibt ihm GOtt auch so herrlichen Sieg wider seine Feinde; auf daß man lerne, alle Sachen mit Gott anfangen, so wird Glück und Heil dabei sein, sonderlich in großen Regimenten, da menschliche Weisheit, Gewalt und Vermögen viel zu schwach ist, daß man der Sachen helfen könnte.

\section{Votum}

Der treue GOtt regiere unsere Obrigkeit, daß sie ferner Rath bei ihm suche, und helfe, daß wir auch zu unserm Beruf den Segen von ihm bitten, und daß Obrigkeit regiere, und die Unterthanen handeln, arbeiten und ordnen, daß es gut, nützlich und selig sei, in Christo Jesu, unserm HErrn, Amen.

\chapter{Das 16. (sonst 15.) Kapitel}

\section{Vorrede}

Die Regenten sollen ihrer weltlichen Hände also warten, daß sie gleichwohl daneben der Ehre Gottes nicht vergessen, und ihren Eifer gegen den Gottesdienst, den zu schützen und zu erhalten, nicht sollen fallen lassen. Sollen auch anordnen und darüber halten, daß die Kirchendiener in Verrichtung der Gottesdienste fleißig und treu sein. Darnach soll der Gottesdienst ein gemeiner Dienst für alle Menschen sein, und sollen diesen mit aller Reverenz und Ehrerbietung in Heiligkeit handeln und verrichten, Jung und Alt, Reich und Arm, geistliche und weltliche Personen, auf daß niemand durch böse Exempel geärgert, und daß Gott die Nachlässigkeit und ungebührlichen Wandel nicht strafen dürfe. Dies erinnert Davids Exempel und Usas Strafe, davon das 16. Kapitel.

\section[Die Bundeslade wird an ihren Ort gebracht]{Die Bundeslade wird an ihren Ort gebracht}
\quellseite{605}

\section{Summarium}

Hie ist eine Historie, wie David die Lade des HErrn läßt gen Jerusalem bringen, an die Stätte, die er bereitet hatte, B. 3. Und nun fort anordnet, daß die Leviten die Lade tragen sollten, B. 12. Denn weil Gott solches zuvor durch Mose geboten hatte, wollte es sich nicht gebühren, daß man anders denn solchem Befehl nachhandeln sollte. Aus diesem Kapitel scheint es, daß Usa kein Levit gewesen und wider den Befehl Gottes die Lade habe angerühret, welches allein den Leviten Amts halber gebühret und derhalben seine Strafe hat tragen müssen.

\section{Beten}

GOtt regiere unsere Obrigkeit, daß sie ferner GOttes Ehre fleißig fördere, und helfe, daß wir ihm auch dienen, damit es ihm gefalle; und er als ein gnädiger, barmherziger GOtt, die heiligen Engel und alle liebe Christen, uns auch dienen an Leib und Seele, wie es unsere geistliche und leibliche Nothdurft erfordert, Amen.

\chapter{Das 17. (sonst 16.) Capitel}

\section{Vorrede}

Wir alle sollen Gott für seine Wohlthaten dankbar sein. Zur rechten Danksagung aber geböret erstlich, daß wir erkennen, daß wir alle gute Gaben von Gott haben. Zum andern, daß wir dieselbigen nicht geringe, sondern hoch und groß achten. Zum dritten, daß wir Gott darum loben und preisen sollen, und nicht allein mit schlechten Worten, sondern auch mit freudigem Gesange. Wie gethan haben Mose, Debora, Hanna, David, Josia, Josaphat, Zacharias, Maria und andere; und solches lehret St. Paulus an die Epheser im 5. Kapitel B. 18. 19., da er sagt: Werdet voll Geistes, und redet unter einander von Psalmen und Lobgesängen und geistlichen Liedern, singet und spielet dem HErrn in euren Herzen, und saget Dank allezeit für alles Gott und dem Vater, in dem Namen unsers HErrn Jesu Christi. Zum vierten, daß wir darum Gott zu Gefallen leben, und uns nach seinen Geboten halten sollen, wie wir auch gemeiniglich geloben, wenn wir unsere Nothdurft von Gott bitten. Und zum fünften, daß wir die Wohlthaten und Güter zu dem Brauch sollen wenden,
\quellseite{606}
dazu sie uns Gott bescheret und gegeben hat. Dies sollen wir uns erinnern bei dem Dankpsalm des Königs David aus dem 17. Kapitel.

\section[Davids Lobgesang und Bestellung des Gottesdienstes]{Davids Lobgesang und Bestellung des Gottesdienstes}

a Diese drei hebräischen Wörter, Cicar, Espar, Assa, heißen nicht allein die Materie, als Brod, Fleisch, Wein, sondern auch das Maß oder Gewicht. Als so ich auf deutsch spräche: Er gab ein Loth Brod, ein Pfund Fleisch, ein Nößel Wein.
\quellseite{607}

\section{Summarium}

Dies ist ein Lied, in welchem das Volk vermahnet wird, daß sie Gottes Güte erkennen und ihm dafür von Herzen danken sollen, daß er sich ihrer so väterlich hat angenommen, seinen Bund mit ihnen gemacht und den rechten Gottesdienst ihnen geordnet hat. Darum vermahnet er, daß sie sich mit Fleiß zu solchem Gottesdienst halten und da Gott loben, seine Wohlthat rühmen und preisen sollen, B. 8. Denn sonst ist nichts, damit wir Gott können dankbar sein. Er befiehlt aber auch, daß man mit Fleiß wolle bitten, daß er sie also in seinem Wort und rechten Gottesdienst erhalten wolle, B. 34. Denn solches ist fürnehmlich vonnöthen, weil der Teufel dem Wort und rechtem Gottesdienst so feind ist.

\section{Votum}

Wir danken Gott für alle seine Wohlthaten, geistliche und leibliche, und der ganzen Christenheit erzeiget; der verleihe, daß wir derer recht gebrauchen zu seiner Ehre, auf daß er uns ferner aus Gnaden segne und behüte, und erhalte uns bei seinem Wort und rechten Gottesdiensten, Amen.

\chapter{Das 18. (sonst 17.) Kapitel}

\section{Vorrede}

Gottes Wort und gottesfürchtiger Leute Exempel und Fürbilder sollen unsers Lebens Regel, Wegweiser und Zuchtmeister sein. Von Gottes Wort lehret uns solches St. Petrus 2 Epistel 1. Kapitel V. 19. und sagt: Daß man auf das prophetische und also göttliche Wort achten soll, und im ganzen Leben und Wandel darauf sehen, als auf ein Licht, das da scheinet und uns fürleuchtet in den Finsternissen dieser Welt. Dann von Exempel der Heiligen unterrichtet St. Paulus an die Philipper 3. Kapitel V. 17. und spricht: Folget mir, lieben Brüder (verstehet auch Schwestern) und sehet auf die, die also wandeln, wie ihr uns habt zum Fürbilde. Wohl nun uns allen, so wir nach diesem Bericht leben. Denn wer Gottes Wort und christlichen Tugendexempeln folget, der hat die Verheißung und Trost, daß es Gott belohnen und daß er seine Zusagungen an seinen Gehorsamen erfüllen will. Derwegen so wollen wir christliche Tugenden aus dem göttlichen Wort und vom Könige David lernen, von welchem das 18. Kapitel.

\section[Der Tempelbau Davids wird eingestellt; und der Messias]{Der Tempelbau Davids wird eingestellt; und der Messias verheißen}
\quellseite{608}

a) (Ich) daß ich nicht meine Person, sondern meines Blutes Nachkommen, scilicet in futurum et longinquum. Der ein solcher Mensch sein wird, der in der Höhe Gott der Herr ist, Ps. 89, 7. Wer kann Gott gleich sein unter den Kindern Gottes? Er ist auch Gottes Kind, aber wohl über andere Gottes Kinder, als der selbst auch Gott ist.

\section{Summarium}

Suche dasselbe im zweiten Buch Samuelis, am 7. Capitel.

\section[Notum]{\unsicher{Notum}}

Der liebe GOtt gebe uns auch gottseelige Herzen, und helfe, daß wir den Tugenden der Heiligen GOttes folgen, ihm dankbar sein, fleißig beten und andere gute Werke üben, damit wir ihnen auch in das ewige Leben folgen, um JEsu Christi, unsers HErrn willen, Amen.

\chapter{Das 19. (sonst 18.) Capitel}

\section{Vorrede}

Selig und überselig ist das Regiment, dem solche Obrigkeit vorstehet, der Gott beistehet und derer Rathschläge er regieret, fortsetzet und zu gutem Ende führet. Denn solche sind HErr Gottes Leute, unter derer Regierung hat die Kirche ihre Herberge, und die Gerechtigkeit ihren Platz; da blühet der Unterthanen Nahrung, da sind die Stadt- und Landwächter die heiligen Engel, da hat ihr Paradies die liebe Jugend, da haben ihren Zulauf, Rettung und Platz die Armen und Verlassenen. Da ist der Friede eine Krone der Unterthanen. Derwegen laßt uns heute abermal Gott herzlich danken für unsere Obrigkeit, die er bisher gegeben, erhalten und gesegnet hat. Ach der helfe, daß bis an jüngsten Tag dies Land und diese Stadt ferner regieret werde von Obrigkeit, die Gott liebe und ehre, welche die Unterthanen ihr treulich lasse befohlen sein, und die christlich, gerecht und friedlich alle ihr Amt verrichte, Amen. Darauf höret an das 19. Capitel.
\quellseite{609}

\section[David bestellet, nach glücklich geführten Kriegen, Amtleute]{David bestellet, nach glücklich geführten Kriegen, Amtleute}

a Segnen) Glück wünschen.

\section{Summarium}

Suche dasselbe im zweiten Buch Samuelis, am 8. Kapitel.

\section{Datum}

O Gott, aller frommen Regenten Rathgeber und Schußherr, lasse ihm auch allzeit unsere liebe Obrigkeit befohlen sein, und verleihe ihr ferner glückselige Regierung, auch Sieg wider alle ihre Feinde, und den Unterthanen fromme Herzen und willige Hände, zum Gehorsam und Nothdurft, gegen der Obrigkeit, Amen.

\chapter{Das 20. (sonst 19.) Kapitel}

\section{Vorrede}

Auf daß die Sünden nicht ungestraft bleiben, so muß der Mensch selber oft Ursache zum Strafen geben. Derwegen, so wollen wir heute unsere Sünden erkennen und in wahrem Glauben herzlich bitten, Gott wolle nicht mit uns handeln nach unsern Sünden, und nicht vergelten nach unsern Missethat, Ps. 103, 10. Wolle auch nicht nach unserm Verdienst uns in seinem Zorn strafen, sondern wolle nach seiner Güte und großen Barmherzigkeit uns gnädig sein, und uns regieren, daß wir nach dem Fleische nicht leben, sondern durch den heil. Geist des Fleisches Geschäfte tödten, Röm. 8, 13., damit wir ihn mit mehr und größern Sünden ferner zu strafen nicht zwingen und bringen. Denn die Uebertretung und Ungehorsamkeit empfähet doch endlich ihren rechten Lohn, Ebr. 2, 12. Und darauf höret an das 20. Kapitel.

\section[David rächet sich an den höhnischen Ammonitern]{David rächet sich an den höhnischen Ammonitern}
\quellseite{610}

a Wie man ein Ding kehret, hinten und vorn besiehet, das man wil eigentlich erkunden.

\section{Summarium}

Suche dasselbe im zweiten Buch Samuelis, am 10. Kapitel.

\section{Votum}

GOtt sei uns gnädig, vergebe uns alle unsere Sünde, und strafe uns ja nicht, wie wir es auch mit unsern großen Sünden vielfältig verdienet haben; und helfe, daß wir allzeit üben eine gute Ritterschaft, behalten den Glauben und ein gut Gewissen, Amen.

\chapter{Das 21. (sonst 20.) Kapitel}

\section{Vorrede}

St. Paulus sagt 1 Cor. 10, 13.: Gott ist getreu, der euch nicht lässet versuchen über euer Vermögen, sondern macht, daß die Versuchung so ein Ende gewinne, daß ihrs könnet ertragen. Dieser Spruch lehret, daß der Teufel, Tod und alle Welt viel zu schwach dazu sind, einen Christen anzufechten und ein Leid zuzufügen, wenn es Gott nicht nachläßt und verhänget. Daraus folget, weil Gott über alles mächtig ist und ohne seinen Willen nichts geschehen kann; sein Wille aber allzeit gut ist: so kann uns dasjenige, so uns nach seinem Willen begegnet, nicht schädlich sein, es dünke die Vernunft so böse es gleich immer wolle, und thue Fleisch und Blut so weh es immer kann. Ferner lehret St. Paulus: Daß uns Gott nur lasse versuchen, aber nicht verderben; es soll nur ein Versuchen, ein Probieren und Bewähren sein, daß wir also geübet werden im Glauben, Liebe, Hoffnung, Geduld, Gehorsam und in dergleichen mehr göttlichen
\quellseite{611}
Stücken. Auch tröstet St. Paulus, daß Gott der Versuchung ein Ende will machen, darum, daß wir stets sein in seiner väterlichen Fürsorge. Er sorget für uns, hütet und wacht, es stehet alles in seiner Macht, und kann die Versuchung, daß Anliegen und Kreuz so stark, geschwind, mächtig und groß nicht sein, Gott ist stärker, mächtiger und größer, und weiß dem allen ein gut und selig Ende zu machen. Das alles bestätigt der König David mit seinem Exempel, von dem das 21. Kapitel.

\section[Davids Krieg und Sieg wider die Ammoniter und Philister]{Davids Krieg und Sieg wider die Ammoniter und Philister}

\section{Summarium}

Hie ist eine Historia, wie Gott dem frommen David, ob er ihn wohl mit vielen Kriegen und Widerwärtigkeiten hat betrüben lassen, dennoch immer heraus geholfen und ihm gnädigen Sieg verliehen, nach dem Spruch Ps. 34, 20.: Der Gerechte muß viel leiden, aber der Herr hilft ihm aus dem allen.

\section{Votum}

GOtt, unserm HErrn, sey Lob und Preis, der zu aller Zeit den Seinen unter dem Kreuze beigestanden ist; der helfe heute auch allen, die in Noth sind, und stehe uns allen bei in der letzten Noth, welche ist der bittere Tod, und erhalte uns zum ewigen Leben, Amen.

\chapter{Das 22. (sonst 21.) Kapitel}

\section{Vorrede}

Im andern Buch von den Maccabäern im 6. Kapitel B. 31. stehet vom heiligen Eleasar, daß seine Beständigkeit im Bekenntniß göttlicher Lehre aufgeschrieben sei, daß sie nicht allein die Jungen, sondern Jedermann zur Tugend ermahnen solle: Also ist die Geschichte von Davids Sünden und derer Strafe fürgeschrieben, daß Gottesfurcht daraus lernen soll Jedermann. Denn Gottes Wort ist geschrieben uns zur Warnung, auf welche das Ende der Welt kommen ist, sagt St. Paulus 1 Cor. 10, 11. Und 2 Tim. 3, 16. spricht er: Alle Schrift von Gott eingegeben ist nütze zur Lehre, zur Strafe, zur Besserung, zur Züchtigung in der Gerechtigkeit, daß ein Mensch Gottes sei vollkommen, zu allem guten Werk geschickt. Denn wie man sich beide gegen Gott und den Menschen halten solle, was man auch thun solle und lassen, solches alles lehret die Schrift der heiligen Biblia, wie mit Exempeln, also mit Sprüchen, durch Befehl, Neigung und warnungsweise. Derwegen höret dazu an das 22. Kapitel.

\section[Dreitägige Pestilenz durch Zählung des Volks verursacht]{Dreitägige Pestilenz durch Zählung des Volks verursacht}
\quellseite{612}

a Thörlich) David bekund die eitle Ehre, daß er wollte sehen, wie das Königreich durch ihn, als durch sein Thun, so herrlich hätte zugenommen. Da siehet du, daß fiducia operis, eitel Ehre wider Gottes Ehre und eine Abgötterei ist.

b Säcken) das ist, mit geringen, groben Kleidern, daraus man Säcke macht, wie bei uns die Kittel und grobe Leinwand.
\quellseite{613}

\section{Summarium}

Suche dasselbe im zweiten Buch Samuelis, am 24. Kapitel.

\section{Votum}

Gott sei uns allen gnädig, und weil wir auch mit unsern Sünden Strafe verdienet haben, so wolle er uns nie hindern, um Jesu Christi, unsers HErrn willen, und wolle uns regieren, daß wir das Böse lassen und das Gute thun allezeit, fleißig und beständig, Amen.

\chapter{Das 23. (sonst 22.) Kapitel}

\section{Vorrede}

Unser HErr Christus sagt Matth. 10, 41.: Wer einen Propheten aufnimmt in eines Propheten Namen, der wird eines Propheten Lohn empfangen. Und wer dieser Geringsten einen nur mit einem Becher kaltes Wassers tränket, in eines Jüngers Namen; wahrlich, ich sage euch, es wird ihm nicht unbelohnet bleiben. Mit diesen Worten lehret der HErr Christus, daß reichlich solle belohnet werden, was man gibt und mittheilet zur Erhaltung des heiligen Predigtamts. Das sollte nun alle Menschen reizen und bewegen, daß sie auf bloßen Knien Christo sollten zutragen, dadurch das Predigtamt gefördert und erhalten wird, weil der Mund und der Grund der Wahrheit verheißt, ja mit einem hohen Eide beteuert, und seine Wahrheit darüber zu Pfande setzet; und der HErr ist, der nicht verarmen kann, daß ers ihnen bezahlen wolle, und zwar mit Propheten- und gerechten Lohn, das ist, daß sie eben solche reiche Belohnung, doch aus Gnaden, bekommen sollen, als die Propheten und Gerechten selbst haben empfangen. Derwegen ein jeder nach seinem Beruf und Vermögen die Kirchen, deren christliche Lehrer und die Gottesdienste versorgen soll helfen, wie auch lehret das 23. Kapitel.

\section[Nothwendige Vorbereitung zum Tempelbau]{Nothwendige Vorbereitung zum Tempelbau}

a Friede) Salomo heißt friedsam oder friedreich.

b Ein Centner ist bei tausend Floren, doch an einem Ort geringer oder mehr \&c. In Gräcia ist 600 Kronen
\quellseite{614}

\section{Summarium}

Hier ist eine schöne Historia, wie David seinem Sohn Salomo das Regiment übergibt, und ihm sonderlich befiehlt, er solle GOtt und sein Wort für Augen haben und demselbigen mit ganzem Ernst nachleben, so werde er glücklich sein, V. 13. Insonderheit aber vermahnet er ihn, den Tempel zu bauen, an welchem David allen Vorrath verschafft hatte, V. 14.; auf daß Salomo sich den Unkosten nicht ließe zurücktreiben. Also sollten zu unsern Zeiten, beide Obrigkeit und Unterthanen, auch thun, die Kirchen vermachen bestellen und versorgen, daß die Nachkommen einen Vorrath finden. Denn es ist zu besorgen, unsere Nachkommen werden sich des Evangelii nicht sonderlich annehmen, noch dasselbige zu fördern sich etwas kosten lassen; weil wir bereits jetzt so kärglich und eigennützig damit umgehen, die wir doch wissen, was für ein Jammer unter dem Papstthum, Leibes und Seelen halber, gewesen ist.

\section{Votum}

Helfe der treue GOtt, daß wir uns nach dieser Lehre gehorsamlich verhalten, und mit Geldsteuern, Testamenten und Gebet zur Kirchen Erhaltung dienen, auf daß bis an jüngsten Tag allhie eine heilige Gemeine sei, die den HErrn Christum recht erkenne und ihn allezeit lobe und preise, Amen.

\chapter{Das 24. (sonst 23.) Kapitel}

\section[Der Leviten Zahl, Ordnung und Amt]{Der Leviten Zahl, Ordnung und Amt}

a. Dreißig) es scheinet hie die ebräische Bibel verfälschet zu sein, denn sonst allenthalben zwanzig geschrieben stehet, ut paulo infra in 24. paragrapho et in Mose.
\quellseite{615}

\section{572}

b Mutatio onerum Mosis. D. i. Aenderung der Lasten Mosis.

\chapter{Das 25. (sonst 24.) Kapitel}

\section[Obersten im Heiligthum und ihre Diener]{Obersten im Heiligthum und ihre Diener}
\quellseite{616}

\chapter{Das 26. (sonst 25.) Kapitel}

\section[Ordnung der heiligen Sänger]{Ordnung der heiligen Sänger}

*Horn zu erheben) ich achte, daß dieser Heman sei Davids Prophet gewesen im königlichen Geschäften, die das Königreich belanget haben, wie er hat streiten und regieren sollen. Denn Horn bedeutet in Regiment und Königreich.
\quellseite{617}

\chapter{Das 27. (sonst 26.) Kapitel}

\section[Bestellung der Thorhüter, Schatzmeister und Richter]{Bestellung der Thorhüter, Schatzmeister und Richter}

a Helden) Denn die Priester mußten zur Zeit des Krieges die \unsicher{Vordersten} im Heer sein mit den Trompeten.
\quellseite{618}

\chapter{Das 28. (sonst 27.) Kapitel}

\section[Davids Feldobersten, Fürsten, Amtleute und Hofräthe]{Davids Feldobersten, Fürsten, Amtleute und Hofräthe}
\quellseite{619}

\chapter{Das 29. (sonst 28.) Kapitel}

\section{Vorrede}

Von dem Nutzen der Furcht Gottes sagt Sirach im 1. Kapitel V. 10. also: Die Furcht des HErrn ist Ehre und Ruhm, Freude und eine schöne Krone. Die Furcht des HErrn macht das Herz fröhlich, und gibt Freude und Wonne ewiglich. Das ist, innwendig macht die Furcht Gottes das Herz fröhlich, macht und des Sohnes Gottes Segen theilhaftig, bringt und gibt gewissen, starken und festen Trost im Tode, und sonst in allen Aengsten und Nöthen; sie rechtfertiget das Herz aus Gnaden durch die zugerechnete Gerechtigkeit des HErrn Christi, und führet wieder an im Herzen und Leben bei den Menschen einen neuen Gehorsam, oder christlichen Wandel, gibt innerlichen Frieden und ewige Seligkeit. Von außen aber gibt die Furcht Gottes immer Glück und Wohlfahrt in allen Sachen, daß es dem Menschen wohl geht, bringet sie zu Ehren, gibt Reichthum und Vollauf, lindert und mindert manchmal viel Strafen und errettet aus leiblicher Gefährlichkeit. Derwegen wie sich jedermann der Furcht Gottes befleißen soll, also sollen die Eltern sonderlich ihre Kinder zur Furcht Gottes ermahnen und ziehen, und sollen beides ja fleißig und treulich thun. Davon höret auch an das 29. Kapitel.

\section[Davids letzter Reichstag, wegen des Tempelbaues angestellt]{Davids letzter Reichstag, wegen des Tempelbaues angestellt}
\quellseite{620}

\section{Summarium}

Sie ist ein sehr fein Exempel, an welchem alle Eltern, sonderlich aber die Fürsten und Herren, lernen sollten, wie sie ihre Kinder unterweisen und wozu sie sie fürnehmlich ziehen und halten sollen, nämlich zur Furcht Gottes, daß sie GOtt erkennen und ihm von Herzen dienen lernen. Wirst du solches thun (spricht David zu Salomo), so wirst du GOtt finden, V. 9. Das ist, GOtt wird mit Segen und Glück bei dir sein. Wirst du ihn aber verlassen, so wird dich GOtt ewig verwerfen. Aber die Welt folget solchem heiligen Exempel nicht; ja sie thut das Widerspiel, darum ist auch so viel Jammers in der Welt.

a Salomo baut aus Gottes Befehl, nach dem Wort Gottes, seinem Vater David zugesagt, 1 Chron. 18, 12. Den selbstgewählten Gottesdienst und Werk mag er nicht.

37
\quellseite{621}

\section{Psalm}

Helfe GOtt, daß die Eltern ihre Kinder in der Furcht GOttes aufziehen, und regiere die Kinder durch seinen guten Geist, daß sie ihm und den Eltern gern und beständig gehorsam sein, auf daß sie in Gottesfurcht leben, auch zeitlichen und ewigen Segen haben und behalten, Amen.

\chapter{Das 30. (sonst 29.) Capitel}

\section{Vorrede}

St. Paulus sagt von den Philippern, in der Epistel an sie geschrieben im 4. Kapitel, V. 14. 18.: Ihr habt wohl gethan, daß ihr euch meines Trübsals angenommen habt. Ihr aber von Philipp(en) wisset, daß von Anfang des Evangelii keine Gemeine mit mir getheilet hat, nach der Rechnung der Ausgabe und Einnahme, denn ihr allein. Denn gen Thessalonich sandtet ihr zu meiner Nothdurft einmal, darnach aber einmal. Ein süßer Geruch, ein angenehm Opfer Gott gefällig. Derwegen dem Predigtamt wohl zu thun und mitzutheilen vergesset nicht, solche eure Opfer werden Gott auch wohl gefallen. So hat Gott auch ein Register, darein er schreibet all dasjenige, was man thut ihm und seinem Predigtamt; und nach dem ein jeglicher Gutes oder Böses gethan wird haben, darnach will er auch einen jeglichen lohnen, wenn er über alles Rechnung halten wird am jüngsten Tage. Von dieser Versorgung des Predigtamts lehret auch das 30. und letzte Capitel.

\section[Steuer zum Bau des Tempels. Salomos und Zadoks]{Steuer zum Bau des Tempels. Salomos und Zadoks Salbung. Davids Abschied}
\quellseite{622}

\section{Summarium}

Das ist ein fein Exempel, wie David und das ganze Volk so willig und reichlich ihre Steuer geben zum Exempel, V. 6. Denn weil auf Erden kein nöthiger noch theurer Schatz ist, denn Gottes Wort und der rechte Gottesdienst, sollte solchem Exempel nach jedermann willig mit seinem Geld und Gut dazu helfen. Aber wie viel sich die Welt das Evangelium kosten lasse, und darauf wage, erfähret man leider täglich. Weil nun David sich solches freiwilligen Opfers von Herzen freuet und Gott von Herzen dafür danket, als für eine sonderliche Gnade, ist leichtlich abzunehmen, was für ein Gefallen Gott an unsern schändlichen Gaben habe, und wie er endlich den Unfleiß und Nachlässigkeit strafen werde, daß niemand mit Ernst darnach trachte, daß man die Kirchen-Aemter sammt den Schulen wohl bestellen, seine gelehrte Leute auferziehen und also Gott seinen rechten Tempel, die christliche Kirche, recht aufbauen möchte.

\section{Votum}

Regiere GOtt Obrigkeit und Zuhörer, daß sie bei dem heiligen Predigtamt thun, was zu desselben Förderung und Erhaltung fruchtbarlich dienet, und belohne solches an beiden durch Vermehrung seiner Erkenntniß, durch Stärkung ihres Glaubens und durch Segen an Leib und Seele, Land und Leuten, Gottes- und Ehre, Amen.

37*
\quellseite{623}

\part{Das zweite Buch der Chronika}

\chapter{Das 1. Kapitel}

\section{Vorrede}

Wer ihm selber wohlgefällt, der mißfällt Gott. Derwegen soll niemand in seinem Beruf, Amt und Nahrung auf seine Weisheit und Geschicklichkeit trauen und bauen, sondern auf Gott soll man allein sehen und von ihm Glück und Segen, oder guten Anfang, fröhlichen Fortgang und nützlichen Ausgang, alle Tage im Namen des HErrn Christi bitten. Dies lehret St. Jacobus im 1. Kapitel seiner Epistel, V. 5., da er also sagt: So aber jemand unter euch Weisheit mangelt (vernehmt, auf daß er in seinem Beruf recht wandele und handele, und darin Glück und Wohlfahrt bekomme und habe), der bitte von Gott, der da gibt einfältiglich jedermann, und rücket es niemand auf; so wird sie ihm gegeben werden. Er bitte aber im Glauben und zweifle nicht. Davon unterrichtet auch Salomo in seinen Sprüchen, und sagt im 2. Kapitel V. 2. also: Laß dein Ohr auf Weisheit Acht haben, und neige dein Herz mit Fleiß dazu. Denn so du mit Fleiß darnach rufest, und darum betest, alsdann wirst du die Furcht des HErrn vernehmen, und Gottes Erkenntniß finden. Eben dies lehret auch Salomo mit seinem eigenen Exempel, davon das 1. Kapitel des andern Buchs der Chronika.

\section[Salomos Weisheit und Reichthum]{Salomos Weisheit und Reichthum}
\quellseite{624}

\section{Summarium}

Suchet dasselbe im ersten Theil des 3. Capitels des ersten Buchs der Könige.

\section{Votum}

Gott, der Herr, regiere alle Obrigkeit, daß sie auch Weisheit und glückselige Regierung fleißig von ihm bitte; und erhöre unser Gebet, daß wir thun für unsere liebe Obrigkeit, daß sie Rath und That zu ihrer Regierung habe, ihm zu Ehren und uns zu geistlicher und leiblicher, zeitlicher und ewiger Wohlfahrt, Amen.

\chapter{Das 2. Capitel}

\section{Vorrede}

Ein Mensch, der sich in wahrem Glauben zu Gott bekehret hat und der Gottseeligkeit sich befleißet, so viel als möglich, der ist und wird recht genennet ein heiliger Tempel Gottes. Denn Gott wohnet in ihm, 2 Cor. 6, 16. Die steinernen Tempel oder Gotteshäuser aber werden heilig genennet, daß sie sind zugeeignet zu heiligen Werken und Gottesdiensten. Von diesem Unterschied prediget der heilige Stephanus in Geschichten der Apostel im 7. Capitel V. 47. 50., da er sagt: Salomo bauete Gott ein Haus; aber der Allerhöchste wohnet nicht in Tempeln, die mit Händen gemacht sind, wie der Prophet spricht: Der Himmel ist mein Stuhl, und die Erde meiner Füße Schemel; was wollt ihr mir denn für ein Haus bauen, spricht der Herr, oder welches ist die Stätte meiner Ruhe? Hat nicht meine Hand das alles gemacht? So helfe Gott, daß wir alle seine Tempel und Wohnungen sein und bleiben, in denen er regiere und wirke, und daß wir sein Lob preisen mit andächtigen Herzen, mit dankbaren Zungen und mit ihm wohlgefälligen Werken. Dieser unser barmherzige Gott helfe auch, daß in unsern Tempeln und Kirchen-Versammlungen allzeit sein Wort recht gelehret, die Sacramente recht ausgetheilet und die Menschen sich recht mit Gebet, Danksagung, Glauben und Gehorsam erzeigen und verhalten. Amen. Darauf höret ferner an das 2. Capitel.

\section[Nothwendige Vorbereitung zum Bau des Tempels]{Nothwendige Vorbereitung zum Bau des Tempels}

b Ebenholz soll ein Holz in Indien sein. \unsicher{Ich hieß vielleicht}, das man jetzt Sandelholz heißt.
\quellseite{625}

\section{Summarium}

Ist eine Historia, daß Salomo dem Befehl Gottes nachkommet und sich zum Bauen schicket. Es ist aber hie sonderlich zu merken, daß er dem Könige zu Tyro schreibet: Es sei nicht darum zu thun, als müßte Gott solches Hauses haben, und könnte es nicht gerathen; so doch der Himmel ihm zu enge und zu geringe sei, zu einer Wohnung, B. 6. Sondern daß er, Salomo und sein Volk, eines solchen Orts dürfen, da sie zusammen kommen, Gottes Wort hören, opfern, beten und anders thun mögen, das ihnen von Gott zu thun befohlen war. Solches war bisher an keinem beständigen gewissen Ort geschehen. Nun fortan aber bestimmet Gott den Ort, und befiehlet, daß zu Jerusalem geschehen sollte; das war eine große herrliche Wohlthat.

\section{Votum}

Wir danken Gott, der uns auch solche Wohltat erzeiget hat, und noch, was zu Seelen und Leibes Nothdurft wir bedürfen, uns reichlich mittheilet; der helfe, daß wir uns in seinen heiligen Tempeln als lebendige Gottes-Tempel verhalten, allda lernen, danken, beten und andere heilige Werke verrichten, zu seinem Lobe und uns zu Nutz, Amen.

\chapter{Das 3. Kapitel}

\section[Vom Bau des Tempels]{Vom Bau des Tempels}

a Morija) auf diesem Berge opferte Abraham seinen Sohn. 1 Mos. 22, 2.

b Man hält, Arnan sei der Jebusiter König gewesen, und bekehret zum Gott Abraham; davon in Commentaren zu reden.
\quellseite{626}

\chapter{Das 4. Capitel}

\section[Von des Tempels zugehörigen Sachen]{Von des Tempels zugehörigen Sachen}
\quellseite{627}

\chapter{Das 5. Kapitel}

\section{Vorrede}

Gleich wie ein menschlicher Leib ohne Seele kein Leben, Sinn noch Bewegung hat: also was ist ein Tempel, Kirche oder Gotteshaus ohne Gottes Wort und rechten Gottesdienst? Derwegen weil wir beides haben, Gotteshäuser und darin das s(w)ligmachende Evangelium, und heilige Ceremonien; so laßt uns auch diese Wohlthaten zu Gottes Ehre recht gebrauchen. Gottes Wort sollen wir mit Andacht hören, fleißig lernen, und fruchtbarlich gebrauchen, zu gottseligem Wandel und nöthigem Trost, im Leben und Sterben. Daneben sollen wir Gott, denn er hats verdient mit seinem und mitgetheilten Gaben; beten sollen wir mit einander, denn unsere Nothdurft erfordert es, zu des Leibes und Seelen Nutz und Gedeihen; die heiligen Sacramente sollen aus gespendet und gebraucht werden, damit wir der heiligen Dreifaltigkeit werden eingeleibt und auf daß unser Glaube werde gestärket und wir also leben mögen, wie es Söhnen und Töchtern Gottes gebühret und wohl aufstehen will. Darauf höret nun an das 5. Kapitel.

\section[Einweihung des Tempels angefangen]{Einweihung des Tempels angefangen}

\section{Summarium}

Ist eine Historia, wie man den Tempel eingeweihet hat, B. 2., das ist, den Gottesdienst darin angefangen habe. Weil aber solches alles geschehen ist nach dem Wort und ausgedrücktem Befehl Gottes, erzeiget sich Gott selbst da, daß das Haus mit einer Wolke erfüllet ward, und niemand drinnen konnte bleiben für der Herrlichkeit des Herrn etc.

Die Papisten haben auch ihre Kirchweihen, aber den Befehl Gottes haben sie nicht. Darum findet man den Teufel eher daselbst, denn unserm Herrn Gott, welchem niemand kann dienen, er folge denn lauter und bloß seinem Wort und Willen. Derhalben sind solche Papisten-Kirchen eitel Mörder-Gruben, sintemal sie nichts anderes darin thun, denn daß sie die armen Gewissen martern, und unzählig viel Seelen morden. Denn da höret man nichts von der Lehre der Gerechtigkeit vor Gott, vom Glauben, vom Gebet etc. Sondern eitel Menschen-Lehre.
\quellseite{628}
Wo aber die regieren, da müssen die Leute Noth an ihrer Seele und Seligkeit leiden. Wie denn Christus selbst klaget, daß die Phariser und Schriftgelehrten aus seinem Hause, daß ein Bethaus allen Völkern sollte sein, eine Mördergrube gemacht haben.

\section{Psalm}

Danket dem HErrn, unserm GOtt, daß wir auch unsere Tempel und rechte Gottesdienste haben; der helfe, daß wir ihn darin mit heiligen Gedanken, Worten und Werken loben und preisen, und wie wir in der Taufe und bei dem Gebrauch des heiligen Abendmahls zugesagt, für Sünden uns hüten und in guten Werken täglich fleißig uns üben, Amen.

\chapter{Das 6. Kapitel}

\section{Vorrede}

Das Beten soll eines jeden Menschen fürnehmstes und bestes Werk sein täglich. Daher sagen die Alten, daß die Hände empor gen Himmel heben (und solches, wenn man betet) sei die beste Arbeit eines Menschen. Wer aber fruchtbarlich beten will, der soll dies allein bitten, was mit Gottes Wort übereinstimmt, und soll glauben, daß Gott die Sünde vergebe und die Sündenstrafe lindern und abwenden wolle, um Christi willen, und daß er geben wolle, was zu seiner Ehre und zu des Menschen Wohlfahrt dienen soll. Davon sagt der heilige Apostel Johannes in seiner 1. Epistel im 5. Kapitel V. 14. also: Das ist die Freudigkeit, die wir haben zu Gott, daß, so wir etwas bitten nach seinem Willen, so höret er uns. Und so wir wissen, daß er uns höret, was wir bitten, so wissen wir, daß wir die Bitte haben, die wir von ihm gebeten haben. Dies betrachten wir auch bei dem 6. Kapitel.

\section[Einweihung des Tempels fortgesetzt]{Einweihung des Tempels fortgesetzt}

a Erwählet) merke hie, wie alles muß aus Gottes Befehl geschehen, auf daß ja niemand aus eigener Andacht Gottesdienst anrichte. Denn Salomo die beide die Stadt Jerusalem und die Person David rühmet, daß sie beide von Gott erwählet sind.
\quellseite{629}

b Geredet) nun laß du Salomo auch dein Wort wahr werden, das du Gott geredet hast, zu wandeln in seinem Gesetz, sonst würde Mühe werden.
\quellseite{630}

\section[- Das 2. Buch der Chronika 6. 7]{- Das 2. Buch der Chronika 6. 7}

\section{Summarium}

Suche dasselbe im ersten Buch der Könige, am 8. Kapitel.

Ist ein Gebet, daß Gott erstlich seine Verheißung, so er David gethan hat, halten und vollstrecken wolle. Siehe B. 14. 15. rc. bis zum Ende.

\section{Votum}

GOtt helfe und regiere uns, daß wir in diesen gefährlichen und betrübten Läufften fleißig beten, und erhöre unser Gebet um JEsu Christi, seines lieben Sohnes willen; und gebe im Kirchen-, weltlichen und Hausregiment, was unsere Nothdurft ist und was zu seines heiligen Namens Lob, Ehre und Preiß gelangen soll, Amen.

\chapter{Das 7. Kapitel}

\section{Vorrede}

Gott, der HErr, sagt Esaia im 65. Kapitel V. 24.: Es soll geschehen, ehe sie rufen, will ich antworten, wenn sie noch reden, will ich hören. Und der HErr Christus spricht Matth. 21, 22.: Alles was ihr bittet im Gebet, so ihr gläubet, so werdet ihrs empfahen. Und der heilige Geist sagt durch St. Paulum Röm. 10, 12.: Es ist aller zumal Ein HErr, reich über alle, die ihn anrufen. Da hören wir, daß Gott, die heilige Dreifaltigkeit, der Gläubigen Gebet erhören will, und will ihnen alles geben, was sie bitten, da es ihm zu Ehren gelangen, und dem, der da betet, und denen, für die er bittet, gut und selig sein soll. Denn Gott erhöret uns nicht allzeit nach unserm Willen, sondern nach unserer Seligkeit, das ist, darnach er erkennet, daß uns gut und selig sein soll. O gelobet und gebenedeiet sei Gott, daß er also seine Vorsichtigkeit und väterliche Treue an uns beweiset und erzeiget. Dieß erinnert auch das 7. Kapitel.

\section[Einweihung des Tempels vollzogen]{Einweihung des Tempels vollzogen}
\quellseite{631}

\section{Summarium}

Suche dasselbe im zweiten Buch der Könige, am 8. und 9. Kapitel. Ist eine Historia, wie der HErr Salomo erschienen, und ihm seine Bitte gewähret habe, B. 12. etc.

\section{Votum}

Der barmherzige GOtt erhöre auch allzeit unser christlich und gläubig Gebet, und helfe, daß wir seine Gaben, geistliche und leibliche, also gebrauchen, auf daß er uns die erhalte und mehre, wie es ihm löblich und uns gut und seliglich sein soll, Amen.

\chapter{Das 8. Kapitel}

\section[Salomos gebauete Städte, Herrschaft, Opfer und Schiffe]{Salomos gebauete Städte, Herrschaft, Opfer und Schiffe}
\quellseite{632}

\chapter{Das 9. Kapitel}

\section{Vorrede}

In Kraft des Werkes der Erlösung und des reichen Verdienstes unsers Herrn Jesu Christi, sind auch wir Heiden zum Evangelio kommen, daß auch wir sollen mit Abraham, Isaak und Jakob im Himmelreich zu Tisch sitzen, Matth. 8, 11., und der Wohlthaten Messia hie im Glauben genießen, dort aber die in Ewigkeit anschauen, 1 Cor. 13, 12. Wie denn von diesem hohen allgemeinen Beruf der Heiden zu Christo auch die Weissagungen des Alten Testaments lehren und zeugen: als daß in dem Samen Abrahams, Isaaks und Jakobs alle Völker auf Erden sollen gesegnet werden. Daß dem Helden Christo, aus dem Stamm Juda geboren, sollen die Völker anhangen. Daß alle Könige ihn anbeten, und alle Heiden ihm dienen sollten. Daß die Heiden würden in seinem Licht wandeln. Daß sie ihre Söhne und Töchter in den Armen und auf den Achseln herzu tragen und sich zu dem neuen Jerusalem versammeln werden, und dem HErrn dienen einträchtiglich. Für diese reichen, wahren, ewigen, himmlischen Güter, durch Christum auch uns, die wir Heiden gewesen, erlanget und erworben, sollen wir dankbar sein, wie St. Paulus vermahnet Col. 1, 11. und sagt: Danksaget dem Vater, der uns tüchtig gemacht hat zum Erbtheil der Heiligen im Lichte; welcher uns errettet hat von der Obrigkeit der Finsterniß, und hat uns versetzet in das Reich seines lieben Sohns, durch welchen wir haben die Erlösung durch sein Blut. Dies erinnert auch der erste Theil des 9. Kapitels.

\section[Salomo von der Königin aus Arabien besucht, hat]{Salomo von der Königin aus Arabien besucht, hat ein großes Einkommen}
\quellseite{633}

\section{Summarium}

Suche dasselbe im ersten Buch der Könige, am 10. Kapitel.

Ist eine Historia, daß Salomo seiner Weisheit halben allenthalben bekannt, bis zu B. 17.

\section[Notum]{\unsicher{Notum}}

Wir danken dem gnädigen und gütigen GOtt herzlich, daß er auch unser, als der Heiden, GOtt ist und daß wir auch Kinder des Allerhöchsten genennet werden; der helfe, daß wir wahre Gliedmaßen des Reiches seines lieben Sohnes, unsers Erlösers, des HErrn JEsu Christi, sein und bleiben zeitlich und ewiglich, Amen.

\section[Der andere Theil des 9. Capitels]{Der andere Theil des 9. Capitels}

\section{Vorrede}

Vom Amt und christlicher heilsamer Regierung der Obrigkeit lehret Gottes Wort also: Moses sagt im 5. Buch im 17. Kapitel V. 19.: Ein König soll lernen fürchten den HErrn, seinen GOtt, daß er halte seine Rechte, daß er darnach thue. Und Esaias im 49. Kapitel V. 23. spricht: Die Könige sollen deine Pfleger, und ihre Fürsten deine Säugammen sein. Und im 32. Kapitel V. 1. sagt er: Fürsten werden herrschen, daß Recht zu handhaben. Also spricht auch St. Paulus Röm. 13, 4.: Die Obrigkeit ist Gottes Dienerin, eine Rächerin zur Strafe über den, der Böses thut. Von dergleichen Eigenschaften und Amtspflichten der Obrigkeit lehret der andere Theil des 9. Capitels.
\quellseite{634}

\section{Beten}

Gott lasse ihm unsere liebe Obrigkeit befohlen sein, und helfe, daß ihre Regierung ihm, auch ihnen selbst und den Unterthanen zu Ehren, Nutz und Wohlthat gelange, in Kraft, Hülfe und Beistand des heil. Geistes, Amen.

\chapter{Das 10. Kapitel}

\section{Vorrede}

Der rechte, größte und Gott angenehmste Gottesdienst ist, die Erkenntniß des göttlichen Namens, Anrufung und Danksagung, auch Geduld und Gehorsam in Leiden und Trübsal. Wer nun seinem Gott treulich und wohl dienen will, der glaube seinem heiligen Wort und bekenne es für jedermann, sollte es auch gleich mit Gefahr Leibes und Lebens geschehen. Denn das will Gott haben, als geschrieben stehet Matth. 10, 32.: Wer mich bekennen wird vor den Menschen, den will ich bekennen vor meinem himmlischen Vater. Und St. Paulus Röm. 10. sagt V. 10.: So man von Herzen glaubet, so wird man gerecht; und so man mit dem Munde bekennet, so wird man selig. Auch soll ein jeder Christ Gott um seine Wohlthaten anrufen, und für empfangene Wohlthaten danken, als denn David lehret Ps. 50, 14. und spricht: Opfere Gott Dank, und bezahle dem Höchsten deine Gelübde. Und rufe mich an in der Noth, so will ich dich erretten, so sollst du mich preisen. Und V. 23.: Wer Dank opfert, der preiset mich. Ueber das sollen alle, die Gott recht dienen wollen, ihr Kreuz in Geduld tragen; und davon lehret die Epistel an die Ebräer am 10. Kapitel V. 36., da gesagt wird: Geduld ist euch noth, auf daß ihr den Willen Gottes thut, und die Verheißung empfahet. So helfe der treue Gott, daß wir in seinem Dienst allzeit rechtschaffen und treu erfunden werden, Amen. Ferner höret an das 10. Kapitel.

\section[Trennung des Reichs, und Abfall der zehn Stämme,]{Trennung des Reichs, und Abfall der zehn Stämme, von Rehabeam verursachet}
\quellseite{635}

\section{Summarium}

Suche dasselbe im zweiten Buch der Könige, am 12. Kapitel.

Ist eine sehr feine Historia, in welcher man siehet etc. Sammt den folgenden zweien Paragraphen.

\section{Gebet}

Gott, unser Herrscher, helfe und verleihe, daß unser Obrigkeit ferner ihre Unterthanen nach seinem heiligen Wort und Willen, väterlich und treulich regiere, und strafe ja nicht an uns der Welt Abgötterei; auch wolle er uns geben, daß wir ihm allzeit dienen in Gerechtigkeit und Heiligkeit, Amen.

\chapter{Das 11. Kapitel}

\section{Vorrede}

In 1 Johannis am 5. Kapitel V. 21. vermahnet der heil. Geist, und sagt: Kindlein, hütet euch vor den Abgöttern. Diesem ernsten Befehl Gottes sollen wir alle folgen. Denn der HErr hasset alle Abgötterei, sagt Sirach im 15. Kapitel V. 13. Ferner vermahnet der heil. Geist Ebr. 5, 14. und spricht: Lasset uns halten an dem Bekenntniß. Diesem ernsten Befehl sollen wir alle auch folgen. Denn verleugnen wir, so wird uns Gott auch verleugnen, sagt Paulus 2 Tim. 2, 12. Beides lehret auch das 11. Kapitel.

\section[Rehabeams Kriegsbereitschaft, Städte, Befestigung, Weiber und Kinder]{Rehabeams Kriegsbereitschaft, Städte, Befestigung, Weiber und Kinder}
\quellseite{636}

\section{Summarium}

Sie ist eine sonderliche feine Historia. Nachdem das Königreich Salomo durch Aufruhr zerrennet, und der König Jerobeam einen neuen Gottesdienst im Königreich Israel hatte angerichtet, lassen die Priester und Leviten, samt den andern gottesfürchtigen Leuten, so im Königreich Israel waren, Haus und Hof dahinten und ziehen gen Jerusalem, da der rechte Gottesdienst war, V. 14. Denn ein Mensch soll sich eher aller Güter, ja auch Leibes und Lebens erwägen, denn daß er sich entweder zur Abgötterei zwingen, oder vom rechten Gottesdienste sollte abdringen lassen.

\section{Votum}

Der treue Gott helfe, daß wir auch bei erkannter Wahrheit beständig bleiben, und alle Abgötterei fliehen bis an unser Ende; und verleihe, daß wir geduldig alles darüber leiden und wagen, wenn es die Noth und unser Gewissen erfordert, Amen.

\chapter{Das 12. Capitel}

\section{Vorrede}

Der HErr unser Gott ist der Armen Schutz, ein Schutz in der Noth, darum hoffen auf dich, die deinen Namen kennen, denn du verlässest nicht, HErr, die dich suchen, Ps. 9, 10. Der HErr verstößet nicht ewiglich, sondern er betrübet wohl und erbarmet sich wieder nach seiner großen Güte; denn er nicht von Herzen die Menschen plaget und betrübet, Klagl. Jer. 3, 31. Denn Gott erbarmet sich der Elenden und bedeckt uns unter dem Schatten seiner Flügel, daß wir erfahren, daß er der HErr sei, an welchem nicht zu Schanden werden, die auf ihn bauen. Wer bist du denn, daß du dich für Menschen fürchtest, die doch sterben, und für Menschen Kindern, die als Heu verzehret werden, und verlässest des HErrn, der dich gemacht hat, Esaia 51, 12. Und der uns ja behütet, wie einen Augapfel, Ps. 17. 8. Dieser Lehre und Trostes Exempel ist das 12. Kapitel.

\section[Rehabeam von Sisak überzogen]{Rehabeam von Sisak überzogen}
\quellseite{637}

\section{Summarium}

Suche dasselbe im ersten Buch der Könige, am 14. Kapitel, und hänge dies Stücklein mit daran.

Sonderlich aber ist zu merken, daß Gott die wohlverdiente Strafe gnädig abwendet, wenn man sich vor ihm demüthiget, das ist, wenn man die Sünde erkennet und lässets ihm leid sein, und hernach sich tröstet, daß Gott gnädig sei, und um solche Gnade und Hülfe bittet.

\section{Dom}

Der treue gütige GOtt leite und führe uns, daß wir alle auch rechtschaffene Buße thun, und um Linderung der verdienten Strafen täglich in wahrem Glauben und herzlichem beten; und helfe und behüte uns auch gnädiglich in diesem der Christenheit elenden und betrübten Zustande, Amen.

\chapter{Das 13. Kapitel}

\section{Vorrede}

Es verlässet Gott gar nicht, die an sein Wort sich halten und seinen Geboten gehorsam sein, sondern hilft ihnen aus in der Noth, und erhält sie wider alle ihre geistlichen und leiblichen Feinde. Davon lehret die heilige Schrift also: Durch Esaiam, den Propheten, spricht Gott im 55. Kapitel B. 2.: Höret mir doch zu, und esset das Gute; so wird eure Seele in Wollust fett werden. Neiget eure Ohren her, und kommt her zu mir; höret, so wird eure Seele leben; denn ich will mit euch einen ewigen Bund machen. Und Salomo im 3. Kapitel B. 22. seiner Sprüche sagt: Halt mein Gesetz und meinen Rath, daß wird deiner Seele Leben sein. Ferner 2 Mos. 23, 22. sagt Gott abermal: Wirst du meines Engels, den ich vor dir hersende, Stimme hören, und thun alles, was ich dir sagen werde, so will ich deiner Feinde Feind und deiner Widersacher Widersacher sein. Desgleichen lehret er im 81. Psalm B. 14. und spricht: Wollte mein Volk mir gehorsam sein und Israel auf meinen Wegen gehen, so wollte ich ihre Feinde bald dämpfen und meine Hand über ihre Widersacher wenden. Diese Lehre laßt uns ja zu unsrem Nutz gebrauchen in gegenwärtigen gefährlichen Zeiten. Davon erinnert auch das 13. Kapitel.
\quellseite{638}

\section[Abias, Königs in Juda, Krieg und Sieg wider]{Abias, Königs in Juda, Krieg und Sieg wider Jerobeam, den König in Israel}

\section[Kap. 14. v. 1]{Kap. 14. v. 1}

\section{Summarium}

Sie ist eine schöne Historia, wie Gott dem Theil, so den rechten Gottesdienst hat, Sieg gibt wider den andern Theil, der abgöttisch war, so doch allewege auf jener Seite zween Mann wi-

38*
\quellseite{639}
der einen waren, V. 17. Also strafet Gott den Jerobeam, und an dem Königreich Israel die Abgötterei. Wiederum beweiset er sich, daß ers mit denen wolle halten, die bei seinem Wort bleiben. Ob sie gleich auch nicht so rein und fromm sind, als sie wohl sollten sein. Denn auf dem andern Theil waren auch Sünder; aber Gott lässet sie seines Worts genießen.

\section{Votum}

Der barmherzige GOtt strafe nicht an uns die bisher begangene greuliche Abgötterei, und regiere uns, daß wir seinem Wort gehorchen, auf daß er uns aus Gnaden auch segne und behüte, und unsern Feinden steure und wehre, Amen.

\chapter{Das 14. Kapitel}

\section{Vorrede}

Im Buch der Weisheit im 3. Kapitel V. 9. wird gesagt: Die Gott vertrauen, die erfahren, daß er treulich hält; und die treu sind in der Liebe, läßt er ihm nicht nehmen. Denn seine Heiligen (oder Gläubigen) sind in Gnaden und Barmherzigkeit, und er hat ein Aufsehen auf seine Auserwählten. Also sagt Salomo in Sprüchen 21. Kapitel V. 22.: Ein Weiser, das ist auch ein Gläubiger, gewinnt die Stadt der Starken, und stürzet ihre Macht durch ihre Sicherheit. Daher sagt St. Johannis in der 1. Epistel im 5. Kapitel V. 4.: Unser Glaube ist der Sieg, der die Welt überwunden hat. Denn freilich ist der Glaube ein König und Herr über Sünde, Tod, Teufel, Welt und alles Unglück, so da mag genennet werden. Derwegen läßt uns Gott bitten, er wolle uns wahren Glauben geben, und helfen, daß wir in allen Nöthen Leibes und der Seele seiner Hülfe uns beständig trösten, und die auch erlangen um Jesu Christ, unsers Herrn willen, Amen. Davon höret auch an das 14. Kapitel.

\section[Asa schaffet die Abgötterei ab, und feget wider]{Asa schaffet die Abgötterei ab, und feget wider die Mohren}

\section{Summarium}

Dies ist eine sonderliche Historia, in welcher ersichtlich ein trefflich Exempel der weltlichen Obrigkeit fürgetragen wird, daß sie keine Abgötterei leiden, sondern sie mit Gewalt abthun und ausrotten soll, V. 3. 5. Solches Werk will Gott mit Frieden und Segen reichlich bezahlen, da er dagegen die Abgötterei mit Krieg und allerlei Unglück strafet. Zum andern ist ein fein Exempel hie eines
\quellseite{640}
sonderlichen Glaubens, daß der König Assa sich den großen Haufen der Feinde nicht schrecken läßt, sondern glaubt gewiß, es sei Gott eben so leicht, den Sieg zu geben wider einen großen Haufen, als wider einen kleinen, V. 11. Solcher Glaube, Vertrauen und Zuversicht auf Gottes Hülfe schlägt die Feinde in die Flucht, daß sie mit ihrer Menge wider den kleinen Haufen nichts mögen ausrichten.

\section{Votum}

GOtt, der Vater, behüte uns ferner für aller Abgötterei, und für der List der Abgötter; und JEsus Christus, der Anfänger und Vollender des Glaubens, stärke unsern Glauben, und verleihe, daß wir durch den Glauben auch alle unsere Feinde überwinden, Amen.

\chapter{Das 15. Kapitel}

\section{Vorrede}

Das Wort Gottes ist lebendig und kräftig, und ist ein Richter der Gedanken und Sinne des Herzens, und ist keine Kreatur vor ihm unsichtbar, wie in der Epistel an die Ebräer im 4. Kapitel V. 12. gesagt wird. Und St. Petrus spricht 1. Epistel 1. Kapitel: Ihr seid wiederum geboren, nicht aus vergänglichem, sondern aus unvergänglichem Samen, nämlich, aus dem lebendigen Wort Gottes, das da ewiglich bleibet. Das ist aber das Wort, welches unter euch verkündiget ist. Allhie hören wir, daß Gottes Wort die Kraft hat, wo man es mit Ernst betrachtet, höret und handelt, daß es nimmer ohne Frucht abgehet; sondern allzeit neuen Verstand, Lust, Freude und Andacht erwecket, rein Herz und Gedanken macht. Denn es sind nicht faule noch todte, sondern kräftige, lebendige Worte. So helfe der treue gütige Gott, daß wir sein heilig Wort allzeit zu unserm geistlichen und leiblichen, zu unserm zeitlichen und ewigen Nutz, zur Besserung und Trost, anhören, lernen und gebrauchen, Amen. Der angezogenen Sprüche und Lehre von Kraft des Wortes Gottes ist ein Exempel das 15. Kapitel.

\section[Assa fährt in seiner Kirchenreformation fort]{Assa fährt in seiner Kirchenreformation fort}
\quellseite{641}

\section{Summarium}

Sie ist eine schöne Predigt, in welcher der Prophet Asarja erstlich warnet, daß der Gottesdienst von Tag zu Tag abnehmen, und die Leute je länger je ärger werden, V. 3. Derhalben vermahnet er, man soll an Gott und seinem Wort fest halten, so werde Gott dem Königreich Juda gnädig sein, V. 7. Solche Predigt läßt ihm der fromme König also zu Herzen gehen, und thut alle Greuel der Abgötterei aus dem Lande und unter andern zerreißet er das Stift seiner Mutter, welches sie dem Abgott Mipßeth gestiftet hatte, V. 8. Solchem Exempel sollen alle Fürsten und Herren folgen, und nicht bei einem Vetter, dort einem Oheim hoffieren und um Menschengunst willen die Abgötterei dulden, und Gottes Wort entweder verfolgen, oder nicht mit Ernst darob halten.

\section{Votum}

Wir danken Gott herzlich, daß unsere Sünden durch sein heilig Wort und Gesetz an uns auch gestraft werden; der helfe, daß wir uns daraus bessern, und damit wir auch allem Unglück Leibes und der Seelen dieser und der zukünftigen Welt entfliehen, Amen.

\chapter{Das 16. Kapitel}

\section{Vorrede}

Der heilige Apostel Paulus sagt 1 Cor. 10, 12.: Wer sich lässet dünken, er stehe, mag wohl zusehen, daß er nicht falle. O eine nöthige Vermahnung. Denn unser Fleisch ist so verderbt, daß es bald kann zum Fall kommen, wo Gott die Hand abzieht. Auch ist der Teufel so listig und gewaltig, daß er es auf allerlei Wege versucht, damit er uns fälle. Dazu so bringets uns die Welt so nahe mit ihren Aergernissen und Bosheit, daß sie uns zu sündigen, ja zu groben Lastern leichtlich bewegen kann, wo Gott nicht uns beistehet und erhält, und wo Christus nicht für uns bittet, daß unser Glaube und Gottseligkeit nicht sinke und abnehme. Daher ist recht gesagt: Große Leute sind gefallen, darum kleine und gemeine, damit sie nicht auch fallen, sich wohl fürsehen mögen. Derwegen fromm sein und fromm bleiben in unserer Macht nicht stehet, und wir könnens ja so bald versehen und so wohl stolpern, als irgend ein anderer. Dies lehret auch das 16. Kapitel.

\section[Assa dreifacher Sündenfall]{Assa dreifacher Sündenfall}
\quellseite{642}

\section{Summarium}

Ist eine Historia, in welcher man siehet, daß der feine und fromme König Asa, der so einen festen Glauben zu Gott hatte, da er wider die Mohren krieget, aus solchem Glauben in einen Unglauben fällt und Hülfe sucht wider den König Israel, nicht bei Gott, wie zuvor, sondern bei dem König der Syrer, B. 3. Da solches der Prophet Hanani strafet, B. 7., wird er zornig über ihn, und wirft ihn in Thurm, B. 10. Aber Gott plaget ihn darum, daß er an Füßen krank wird und stirbt, B. 12. In solcher Krankheit (sagt der Text, B. 12.) suchet er den HErrn nicht, sondern die Aerzte. Die Aerzte suchen, ist nicht unrecht, denn Gott die Arznei den Menschen zu gut, sammt andern Kreaturen erschaffen hat. Aber Gott will zuvor gesucht sein, das ist, er will, daß man in Krankheit und allen Nöthen ihn anrufen und Hülfe bei ihm suchen soll, und alsdann die Creatur oder Arznei brauchen, mit der Zuversicht, Gott werde seinen Segen dazu geben.

Solche Exempel sollen uns reizen, daß wir nicht sicher sein, sondern Gott ohne Unterlaß bitten, daß er uns in seiner Furcht erhalten und mit seinem heil. Geist regieren und für Sünden gnädiglich behüten wolle. Denn wo Sünden sind, da findet sich auch die Strafe.

\section{Votum}

Helfe der treue GOtt, daß wir dieser Lehre fleißig und beständig folgen, im Glauben verharren, und in Glück und Unglück christlich uns verhalten. Er wolle ihm auch alle kranke Christen lassen befohlen sein, zu Hülfe und Trost, um JEsu Christi, seines geliebten Sohnes willen, Amen.

\chapter{Das 17. Kapitel}

\section{Vorrede}

Im 5. Buch Mose am 12. Kapitel B. 2. befiehlt Gott, und sagt der ordentlichen Obrigkeit: Verstöret alle Orte, da die Heiden ihren Göttern gedienet haben, und die Gößen ihrer Götter thut ab und vertilget ihren Namen aus demselben Ort. David, der königliche Prophet, spricht die Obrigkeit im 2 Ps. 10. 11. also an:

Und im Buch der Weisheit am 6. B. 22. lehret Gott die Obrigkeit auch also, und spricht: Wollt ihr nun, ihr Tyrannen (oder Fürsten) im Volk gern Könige und Fürsten sein, so haltet die Weisheit in Ehren, auf daß ihr ewiglich herrschet. Durch das Wort Weisheit aber soll verstanden werden Gottes Wort und das heilige Predigtamt, mit dessen treuen Dienern und Predigern. Von solchem der Obrigkeit Amt lehret auch das schöne Exempel Josaphat, von dem das 17. Kapitel.

\section[Josaphat ist bei seiner Gottesfurcht glücklich]{Josaphat ist bei seiner Gottesfurcht glücklich}
\quellseite{643}

a Vorrath b) nicht allein des Getreides, sondern auch des Zeugs, Waffen oder Rüstung.

\section{Summarium}

Dies ist eine sehr feine Historia, wie der fromme König Josaphat im Anfang seines Regiments sich vor allem besinnet, daß er die Abgötterei ausrotte, V. 7., und rechtschaffene Prediger an alle Orte seines Königreichs verordnet, V. 9. Solches ist das fürtrefflichste und ein recht fürstlich Werk, da Gott sonderlich mit Gnaden will bei sein.

Aber der Teufel läßt die Obrigkeit nicht gern zu solchen guten Werken kommen; darum treibt er den meisten Theil dahin, daß sie die frommen Prediger verfolgen und in ihren Landen und Gebieten die rechte Lehre nicht dulden, wie viele Exempel für Augen sind. Wie es aber ein Ende werde nehmen, wird die Zeit ausweisen.

\section{Votum}

Der allmächtige GOtt steure dem Zerstörer alles Guten, dem Teufel, daß er niemand das Gute hindere; und segne die Regierung unserer lieben Obrigkeit, und aller christlichen Potentaten und Herrschaft, daß dieselbe zu Abschaffung aller Abgötterei diene, und daß Gottes-Wort in aller Welt rein und lauter gepredigt werde, Amen.

\chapter{Das 18. Kapitel}

\section{Vorrede}

Die Alten sagen ja schön: Wie der zeitliche Hunger eine gewisse Anzeigung ist, daß der Mensch frisch und gesund sei: also sei es auch eine gewisse Vermuthung der Seelen Gesundheit, wenn ein Mensch Gottes Wort (welches der Seelen Speise wider den ewigen Hunger ist) lieb und werth hat, daß ihn darnach hungert und dürstet. Wie aber eines in einem lieblichen Lustgarten sich erlustiget und darauß ein schönes wohlriechendes Blümlein abbricht: also sollen wir uns auch in dem herrlichen Lustgarten des göttlichen Worts (darinnen die Blümlein des Lebens, als Herzenskraut, Angelica, Benedicten-Wurzel, Ehrenpreis, himmlischer Balsam, Herrgotts Beerlein, Teufels Abbiß, Heilwurzel, Augen-Trost, himmlisches Wundkraut und dergleichen zu finden sind) fleißig umsehen, heilsame Blumen da abbrechen und damit unsere zur Sünde geneigte und auch matte Herzen, die der leidige Satan erstänken und ersticken will, mit dem Geruch des Lebens zum Leben, 2 Cor. 2, 16., wieder erquicken, auf daß wir in wahrem Glauben und gottseligem Wandel verharren, und mit der Welt nicht aus einer Sünde in die andere fallen und damit Gott und Menschen beleidigen und beschädigen. Dazu hört auch an das 18. Kapitel.

\section[Josaphats Zug mit Ahab wider die Syrer]{Josaphats Zug mit Ahab wider die Syrer}
% Scan page 644: omitted; Bible text/narrative only, no retained apparatus.
\quellseite{645}

\section{Summarium}

Suche dasselbe im ersten Buch der Könige, am 22. Kapitel.

\section{Votum}

Der gütige, barmherzige GOtt erleuchte und regiere alle christliche Obrigkeit, daß sie sich durch falsche Lehrer nicht verführen lasse; und steure noch heute allen, die sein Wort verachten, behüte alle treue Lehrer seines Evangelii, tröste und erlöse alle Gefangenen, und segne uns alle zeitlich und ewig, Amen.

\chapter{Das 19. Kapitel}

\section{Vorrede}

Wer seiner Sünden halber öffentlich oder im geheimen, von der Kanzel oder im Beichtstuhl, gestraft wird, der soll (wes Standes und Wesens er auch ist) die Strafe annehmen, von Sünden abstehen und rechtschaffene Früchte der Buße thun. Denn ein jeder Mensch soll Gott dienen nach seinem Beruf, und soll kein Werk unterlassen, daß der christlichen Kirchen und dem weltlichen Regiment zum Besten gelangen mag. Solchen Gehorsam belohnet Gott also väterlich und reichlich, Matth. 10, 42., daß auch nicht ein Becher kaltes Wassers, der Gottes wegen gegeben wird, unvergolten bleiben soll. Ja, er wird alle Nothdurft erfüllen nach seinem Reichthum, in der Herrlichkeit, in Christo Jesu, Phil. 4, 19., denen allen, die sonderlich Gottes Wort und die rechten Gottesdienste befördern und erhalten helfen. Dies soll uns erinnern das 19. Kapitel.

Josaphat von Jehu bestraft; bestellet die Kirche und das weltliche Regiment.

\section{Summarium}

Dies ist erstlich eine Historia, wie Gott den frommen Josaphat durch Jehu, den Propheten, ernstlich strafen lässet, daß er mit dem gottlosen König Ahab sich vereiniget, und ihm wider die Syrer geholfen hat, B. 2. Denn Gott will, daß man mit denen, so öffentlich Abgötterei trei-
\quellseite{646}
ben, keine Gemeinschaft haben und ihnen gar nichts helfen soll. Denn Gott hat im Sinn, sie zu strafen ihrer Abgötterei halben, weshalben wer zu ihnen steht, sich selbst in die Gefahr gibt. Zum andern ist wiederum ein Exempel hie, dem alle Regenten folgen sollten, daß Josaphat das weltliche Regiment und die Kirchen mit tüchtigen und frommen Leuten bestellet, auf daß niemand unrecht geschehe, B. 6. Sonderlich aber, daß die Leute im Wort recht unterwiesen werden, B. 8. Solches soll sich weltliche Obrigkeit mit Ernst befleißigen, denn Gott ein sonderlich Gefallen daran hat, und gibt reichlich seinen Segen dazu.

\section{Votum}

GOtt, der heilige Geist, verleihe und wirke, daß alle Abgötter heute auch ihre Sünden erkennen; und helfe, daß unsere Obrigkeit die Kirche und Regiment ihr ferner lasse befohlen sein, auf daß sie mit allen ihren Unterthanen GOtt lobe und preise, hie und dort in alle Ewigkeit, Amen.

\chapter{Das 20. Kapitel}

\section{Vorrede}

Das 20. Kapitel des andern Buchs der Chronika ist das fürnehmste Kapitel dieses ganzen Buchs, und ist ein schön Exempel des Spruchs des königlichen Propheten Davids im 50. Psalm, V. 15., da er in Gottes Person also sagt: Rufe mich an in der Noth, so will ich dich erretten, so sollst du mich preisen. Er lehret allda, wir sollen in keiner Noth verzagen, sondern Gott um Hülfe in wahrem Glauben und beständiger Geduld bitten, es sei die Noth so groß, so schwer, so breit sie möge, er könne und wolle uns, zu seiner Zeit, daraus erretten, und also helfen, daß wir ihn genug zu rühmen und zu preisen haben werden. Und sollen aus diesem Kapitel wohl und mit Fleiß behalten werden die zween schönen Lehr- und Trostsprüche, als der erste Spruch ist, B. 12.: Wir wissen nicht, was wir thun sollen, sondern unsere Augen sehen nach dir, Gott. Der andere Spruch ist, B. 20.: Glaubet an den HErrn, euren Gott, so werdet ihr sicher sein, und glaubet seinen Propheten, so werdet ihr Glück haben.

\section[Josaphats Sieg wider die Ammoniter und Moabiter]{Josaphats Sieg wider die Ammoniter und Moabiter}
\quellseite{647}

\section{Summarium}

Ist eine schöne Historia, wie der fromme König Josaphat in der Noth sich an den rechten Trost hält und zu Gott um Hülfe schreyet, ihn seiner Zusagung und herrlichen Wohlthat erinnert und bittet, er wolle jetzt in dieser Noth auch dermaßen helfen, V. 6. Solches Gebet erhöret Gott, und schaffet wunderbarer Weise Hülfe aus aller Noth, uns zum Exempel, daß wir auch, wie Josaphat hie heißt, an den HErrn unsern Gott glauben sollen, so werden wir Schutz finden und Glück haben.
\quellseite{648}

\section{Summarium}

Ist erstlich ein sonderlich Exempel eines großen Glaubens, daß Josaphat an des Propheten Jehasiel Wort so gar nicht zweifelt, daß er sein Volk nicht anders zur Schlacht rüstet, denn als sollten sie gen Jerusalem im Tempel gehen, nicht in Streit ziehen, heißt sie singen, Gott danken und seine Güte preisen, als wären die Feinde schon geschlagen und der Sieg erobert, B. 27.

Zum andern ist hie ein Exempel, an welchem wir auch sollten Gott lernen trauen, sintemal er so leicht und unversehens, in der höchsten Noth kann Hülfe und Rath schaffen, wie er hie thut mit den Feinden Josaphat, B. 24.

Am letzten ist wieder ein Exempel, weil Josaphat sich mit dem gottlosen König Ahazia vereiniget, hat er kein Glück, B. 35. Denn Gott will weder Glück noch Heil denen widerfahren lassen, so sich an die Gottlosen, so in öffentlicher Abgötterei leben, hängen, und mit ihnen zu thun haben.

\section[Notum]{\unsicher{Notum}}

Der treue GOtt gebe und erhalte uns den Geist des Gebets, erhöre auch unser Seufzen, Rufen und Schreien in unserer Noth; und regiere uns, daß wir seiner Hülfe und Wohlthaten recht und wohl gebrauchen, ihm zu Lob, uns zu Nutz und unsern Nächsten zur Wohlfahrt, Amen.

\chapter{Das 21. Kapitel}

\section{Vorrede}

Wie der Teufel von Anfang der Welt den Menschen mit allem Fleiß nachgestellet hat, auf daß er sie fällete, ihnen zu zeitlichem und ewigem Schaden, und Gott, dem er spinnfeind ist, zu Unehren: also gehet er noch alle Augenblicke herum wie ein brüllender Löwe, und suchet, welchen er verschlinge, 1 Pet. 5, 8. Gleich wie auch Gott, so lange die Welt stehet, die Sünden an allen Unbußfertigen gestraft hat, zu Lob seiner Gerechtigkeit und zur Strafe der Gottlosen: also will er noch heute strafen alle, die da sündigen und nicht wahre Buße thun. Derwegen laßt uns wachen und beten, auf daß uns der Teufel nicht in Versuchung führe, und damit wir Gottes Zorn, Ruthen und Peitschen nicht verdienen, leiden und mit Noth und Weh fühlen dürfen. Dies lehret auch das 21. Kapitel.

\section[Jorams Brudermord, Abgötterei, Strafe und elender Tod]{Jorams Brudermord, Abgötterei, Strafe und elender Tod}
\quellseite{649}

\section{Summarium}

Dies ist eine Historia, da man erstlich siehet, daß der fromme König Josaphat einem bösen gottlosen Sohn, B. 1., welcher nicht allein allerlei Abgötterei anrichtet, sondern auch seine Brüder alle, dazu auch etliche Obersten in Israel erwürget, B. 4. Darum plaget ihn Gott auch erstlich mit Krieg, daß nicht allein die Unterthanen, Edom und Libna von ihm abfallen, B. 8., sondern auch die Araber und die Philister verwüsten und plündern das Land, B. 16. Auch wird er gestraft am Leib mit heftigen Krankheiten, B. 19. Auf daß jedermann sehe, was die zu gewarten haben, so gottlos sind, und mit Abgötterei umgehen, und wir uns desto fleißiger zu Gottes Wort und dem rechten Gottesdienst halten.

\section{Votum}

Der barmherzige GOtt, unser lieber Vater, regiere uns, daß wir uns den Teufel nicht verführen lassen; und strafe uns nicht, wie wir es alle schon wohl verdienet haben, und steure allem Feinden seiner Christenheit, geistlichen und leiblichen, Amen.

\chapter{Das 22. Kapitel}

\section{Vorrede}

Wer das 22. Kapitel des andern Buchs der Chronika ablesen höret, derselbe soll herzlich erschrecken, daß der zornige und gerechte Gott die zwei königlichen Geschlechter, als Salomonis und Ahabs, so schrecklich gestraft und vertilget hat, und soll ihm ein jedes, reich und arm, edel und unedel, Ursache daraus nehmen, daß es Gott fürchte und die Sünde meide und meide; und soll ihm diese Rechnung machen: Hat Gott nicht verschonet der Geschlechter Salomonis und Ahabs, sondern hat sie der Abgötterei und anderer Sünden wegen gar ausgerottet: so wird er auch anderer nicht verschonen, die mit dergleichen Sünden Strafe verdienet haben; und soll ein jeder also leben, lehren, regieren, handeln, erben und arbeiten, daß er nicht auch zum Verderben und Untergange seines Geschlechts diene und helfe mit seinen Sünden, Lastern und Missethaten. Darauf höret an das gemeldete 22. Kapitel.

\section[Ahazias Untergang. Athalias Mordthat]{Ahazias Untergang. Athalias Mordthat}
\quellseite{650}

\section{Summarium}

Suche dasselbe im andern Buch der Könige, am 9. Kapitel.

Ist eine Historia, wie der gottlose König in Israel, mit Namen Joram etc.

\section{Gebet}

Gott, unser lieber Vater, der es allein thun kann, wende gnädiglich von uns ab die verdienten Strafen; und regiere die Obrigkeit und die Unterthanen, daß sie Abgötterei und alle Sünden fliehen und meiden, und nach Gottes Wort regieren und leben, in Kraft des heiligen Geistes, Amen.

\chapter{Das 23. Kapitel}

\section{Vorrede}

In Sprüchen Salomonis im 24. Kapitel B. 21. sagt der heilige Geist: Fürchte den HErrn und den König, und menge dich nicht unter die Aufrührischen. Denn ihr Unfall wird plötzlich entstehen. Und im 37. Psalm B. 39. 40. spricht David: Der HErr hilft den Gerechten, der ist ihre Stärke in der Noth; und der HErr wird ihnen beistehen und wird sie erretten. Er wird sie von den Gottlosen erretten und ihnen helfen, denn sie trauen auf ihn. Und im Buch der Weisheit im 6. B. 7. stehet: Die Gewaltigen werden gewaltiglich gestraft werden. Dieser Sprüche Exempel ist das 23. Kapitel.

\section[Jojada salbet Joas zum Könige. Athalia wird getödtet;]{Jojada salbet Joas zum Könige. Athalia wird getödtet; Baal zerstöret}
\quellseite{651}

a Zeugniß sein ist dem Könige beides die Krone und das Buch gegeben, auf daß er nicht allein mächtig, sondern auch weise sein sollte, oder (wie man es reden mag) Gottes Wort und Recht wissen. So machet man jetzt Könige mit einem Schwert und Buch \&c.

\section{Summarium}

Suche dasselbe im andern Buch der Könige, am 13. Kapitel.

Ist eine merkliche Historia, wie der Priester Jojada \&c.

\section{Votum}

Der starke und gerechte GOTT steure heute auch allen Tyrannen und Abgöttern, um seines heiligen Namens Ehre willen; und erhalte uns sein Wort und Gottesdienst, ihm auch zu Ehren, uns und unsern Nachkommen zu täglicher Buße, zu täglichem Trost, und zur ewigen Seligkeit, um JEsu Christi, seines lieben Sohnes und unsers HErrn willen, Amen.

\chapter{Das 24. Kapitel}

\section{Vorrede}

Der heilige Prophet Jeremias spricht im 10. Kapitel V. 23.: Ich weiß, HErr, daß des Menschen Thun stehet nicht in seiner Gewalt, und stehet in Niemandes Macht, wie er wandele oder seinen Gang richte. Das ist so viel gesagt: Wenn ein Mensch durch seine Sünden den heiligen Geist ausgejaget und verstoßen, so wird aus dem größten Heiligen und frömmsten Christen ein großer Sünder. Und in seinen Sprüchen im 29. Kapitel V. 18. sagt Salomo: Wenn die Weissagung aus ist, so wird das Volk wild und wüst. Das ist so viel gesagt: Ohne Gottes Wort, welches lehret, wie man recht glauben und christlich leben soll, kann der Mensch nichts anders thun, denn Abgötterei und seinen Willen nach fleisches Lust in allerlei Sünden treiben. Derwegen, ach laßt uns züchtig, gerecht und gottselig leben, Tit. 3, 12., daß wir Tempel und Wohnungen des heiligen Geistes bleiben. Laßt auch Gottes Wort allein unsern Zuchtmeister und Wegweiser sein, in Thun und Lassen, nach unserm Beruf und Stande, daß wir Gott recht dienen unser Lebelang, in Heiligkeit und Gerechtigkeit, die ihm gefällig ist, Luc. 1, 75. Dies soll uns auch erinnern das 24. Kapitel.

\section[Das übrige Capitel, Abgötterei, Strafe und Tod]{Das übrige Capitel, Abgötterei, Strafe und Tod}
\quellseite{652}

\section{Summarium}

Suche dasselbe im zweiten Buch der Könige, am 12. Kapitel.

Ist eine feine Historia, wie Jojada sammt dem Könige Joas ordnen etc.

\section[Der andere Theil des 24. Kapitels]{Der andere Theil des 24. Kapitels}

Ist eine feine Historia, in welcher man siehet, wie großen Fürsten und Herren sonderlich viel daran gelegen ist, was sie für Prediger und Lehrer haben. So lange der Hohepriester Jojada lebte, that Joas recht und hielt sich wohl. Da er aber starb, wandte es sich mit Joas; denn er folgte seinem gottlosen Räthen und ließ sich von ihnen verführen, daß er abgöttisch ward und den Propheten Zachariam, des Hohenpriesters Jojada Sohn, der ihn um solches strafet, steinigen ließ, V. 21. Solches, spricht der Prophet, wird der HErr sehen und strafen, V. 22. Denn Joas ward bald hernach mit Krieg überfallen. Und neben dem, daß ihn Gott plaget mit großen Krankheiten, wird er endlich noch erstochen, V. 25. Zur Strafe, daß er den HErrn verlassen und den Propheten Zachariam unschuldig hatte erwürgen lassen. Zum Exempel allen Regenten, daß sie dergleichen nicht auch thun, oder dergleichen Strafe gewiß gewarten.
\quellseite{653}

\section{Gebet}

Gott, der heilige Geist, regiere und führe uns, daß wir in Gottes Furcht beständig bleiben, für Abgötterei uns hüten und seinem Wort allezeit gehorsam sein; und behüte uns gnädig für zeitlichem und ewigem Verderben und Untergang. Amen.

\chapter{Das 25. Kapitel}

\section{Vorrede}

Salomo sagt in seinen Sprüchen im 24. Kapitel B. 16.: Ein Gerechter fällt sieben mal. Und unser HErr Jesus Christus spricht Luc. 13, 3.: So ihr euch nicht bessert, werdet ihr alle auch also umkommen. Dervwegen, weil der, so heute stehet, morgen fallen kann, und der, so jetzt fromm ist, gar bald sündigen kann, und weil Gott die Sünden nicht will ungestraft lassen; so sollen wir vorsichtig wandeln und uns an den Stab göttlichen Wortes halten, auf daß, weil wir in dieser Welt, als auf einem glatten Eis gehen, nicht fallen; und da ja jemand strauchelt oder einen Stolper thät, der nehme die treuen Bußpredigten an, und gebrauche derer recht, der erkenne, daß er übel und unrecht gethan, der bitte, Gott wolle ihm solche Sünde um des HErrn Christi Verdienstes willen vergeben; der hüte sich für vorgethanem und allem andern Sünden, der lebe christlich und ehrbar, der bete auch, Gott, der heilige Geist, wolle ihm helfen, daß er in Gottesfurcht beständig lebe, damit Gottes Zorn nicht auch wie Wasser auf uns ausgegossen werde, und uns an Leib und Seele, Gut und Ehre zeitliche und ewige Strafe nicht folgen. Dies soll uns auch erinnern das 25. Kapitel.

\section[Regierung Amazias, des Königs in Juda]{Regierung Amazias, des Königs in Juda}
\quellseite{654}

\section{Summarium}

Dies ist eine sonderlich feine Historia, in welcher erstlich zu merken ist, wie Gott dem König Amazia warnen läßt, daß er die hundert tausend Mann, die ihm aus dem Königreich Israel zu Hülfe gezogen waren, wieder soll heim schicken, V. 10. Darum, daß Gott kein Glück geben will, wo man Hülfe bei denen sucht, so gottlos sind, unrechte Lehre und falschen Gottesdienst haben und vertheidigen. Um solcher Ursachen willen ist zu besorgen, daß wir bisher wider den Türken so wenig haben ausgerichtet.

Zum andern ist hie wieder ein schrecklich Exempel, daß uns zur Furcht Gottes halten soll. Nachdem der König Amazia so einen herrlichen Sieg durch Gottes Hülfe erlanget hatte, fällt er von Gott ab, nimmt die Götter der Edomiter (welche er geschlagen hatte), und stellet sie ihm zu Göttern, und betet an für ihnen, V. 14. Und da er von einem Propheten gestraft wird, verdroß es ihn, und dräuet dem Propheten, V. 16. Solche Sünde machts bald aus mit seinem Regiment.
\quellseite{655}

\section{Summarium}

Sie ist eine Historia, in welcher man siehet, ob wohl Amazia eine redliche Ursache hatte wider den König Israel, darum, daß sein Kriegsvolk ohne alle Ursache drei tausend erschlagen, und etliche Städte geplündert; dennoch hat er kein Glück; Ursache war diese, daß er sich an der Heiden Abgötterei gehangen hatte. Um solcher Sünde willen verlieret er die Schlacht, Jerusalem wird geplündert und zerrissen, B. 23., und er, der König Amazia, wird endlich erstochen, B. 28. Solchen Lohn sollen alle die gewarten, so Abgötterei und falscher Lehre folgen und sie handhaben.

\section{Gebet}

Helfe GOtt, daß wir uns für gottloser Leute Gesellschaft und Gemeinschaft hüten; verleihe auch, daß wir uns aus den Strafpredigten bessern und in der Frömmigkeit verharren, und lindere die wohlverdienten Strafen, um JEsu Christi, seines lieben Sohnes willen, Amen.

\chapter{Das 26. Kapitel}

\section{Vorrede}

Obschon ein Mensch eine Zeitlang gottesfürchtig und glückselig ist, und thut in seinem Beruf, was Gott gefällt und den Leuten nützlich ist; dennoch kann er fallen, sündigen, unrecht thun und in zeitlichen und ewigen Jammer kommen, wenn er aus Fürwitz oder Frevel verführet durch sein Fleisch und Blut, oder durch die Welt, oder durch den Teufel wider Gottes Wort und Gebot, in seinem Stande thut und handelt. Dieses sollen wir wohl betrachten und täglich also beten: O HErr, bestätige das gute Werk, das du in uns gewirket hast, und die Werke deiner Hände wollest du ja nicht lassen, Ps. 138, 8. Und wollest solches thun um der Ehre deines heiligen Namens willen, und daß es zu unseres Nächsten Nutz gelange, und uns zu zeitlicher und ewiger Wohlfahrt diene, um des Verdienstes JEsu Christi, deines geliebten Sohnes und unsers HErrn und Heilandes willen, Amen. Das lehret auch das 26. Kapitel.

\section[Usia räuchert und wird aussätzig]{Usia räuchert und wird aussätzig}

a (Gesichte) Gesichte sind Prophezeiungen. Er will sagen, Sacharja sei gelehrt gewesen in Mose und den Propheten, Samuel, David, Gad und dergleichen.
\quellseite{656}

\section{Summarium}

Dies ist erstlich wieder ein Exempel, daß man siehet, wie viel einem Regenten daran gelegen sei, daß er einen rechten Prediger und Lehrer habe. Denn hie stehet man, so lange Zacharias lebet, der Lehrer in den Gesichten Gottes, suchet der König Gott und gibt ihm Glück, daß ihm in allen Dingen wohl gelinget, B. 5. Da aber Zacharias stirbt, erhebet sich Usia seiner Macht und Glückes, verachtet die Priester, und will beides, König und Hohepriester sein, und selbst auf dem Altar opfern, B. 16. Um solcher Sünde willen strafet ihn Gott mit dem Aussatz sein Lebenlang, B. 19. Denn Könige und Fürsten sollen eben sowohl unter Gottes Wort sein und demselbigen gehorchen, als andere Leute; thun sie es nicht, sollen sie eben sowohl ihre Strafe darum leiden, als andere.

\section{Gebet}

GOtt, unser lieber Vater, erbarme dich unser (denn wir bisher seinem Worte auch nicht aller Dinge gefolget haben); und verleihe, daß wir nun ferner seine gehorsamen Kinder sein und bleiben, und von ihm an Leib und Seele gesegnet werden, und in unserm Beruf Glück und Wohlfahrt behalten, Amen.

\chapter{Das 27. Kapitel}

\section{Vorrede}

David sagt im 103. Psalm.: Wie sich ein Vater über Kinder erbarmet, so erbarmet sich der HErr über die, so ihn fürchten. Und B. 17.: Die Gnade des HErrn währet von Ewigkeit zu Ewigkeit, über die, so ihn fürchten. Und im andern Buch der Könige im 17. Capitel B. 39. spricht der heilige Geist: Fürchtet den HErrn, euren Gott, der wird euch erretten von allen euren Feinden. Freilich beschützt Gott alle, die ihn fürchten, wider alle Teufel und wider die ganze Welt, der er mächtig und gewaltig ist und ihnen zu gebieten und zu verbieten hat, und ein solcher HErr ist, der alle Dinge übersiehet und überhöret. Derwegen, wem er wohl ist, dem kann niemand übel; wer ihn zum Freunde hat, dem kann kein Feind schaden; wem er liebet, den kann kein Haß beschädigen. Dieser Sprüche Exempel ist das 27. Capitel.

\section[Jotham macht die Ammoniter zinsbar]{Jotham macht die Ammoniter zinsbar}
\quellseite{657}

\section{Summarium}

Ist eine Historia, daß Jotham, der König Juda, viele Siege und Glück hat, um der Ursache willen, daß er (wie der Text sagt) seine Wege richtet für Gott, seinem Herrn, das ist, er lebet in Gottesfurcht und nach dem Wort Gottes. Zum Exempel aller Obrigkeit: so sie glückselige Regierung begehrt, daß sie sich gottselig halten und nach Gottes Wort in der Furcht Gottes leben soll; wo nicht, so sollen sie alles Unglücks gewarten.

\section{Votum}

Der treue GOtt gebe der Obrigkeit und Unterthanen einerlei Herz, und regiere sie durch seinen Geist und Wort, daß sie ihn fürchten ihr Lebenlang, auf daß es ihnen und ihren Kindern wohl gehe und kein Unglück über sie komme, Amen.

\chapter{Das 28. Kapitel}

\section{Vorrede}

Gottes Wort lehrets und die allgemeine Erfahrung bezeugets, daß Gott fromme Leute segnet und erhält; und dagegen, daß Gott unbußfertige Sünder strafet und verderbet. Derwegen laßt uns alles sündliche Wesen fliehen, und aller christlichen Tugenden befleißigen, mit vorgehendem Gebet um Hülfe und Regierung des heiligen Geistes; so wird unsere Arbeit, Regieren, Lehren, Handeln und Obern nicht eine vergebene, unnütze Mühe sein, sondern wird von Gott gefördert und aus Gnaden mit Glück und Fortgang belohnet werden. Denn also lautet davon die gnädige Verheißung Gottes, Esaiæ 65, 23.: Meine Auserwählten sollen nicht umsonst arbeiten. Und St. Paulus sagt von denen, die ihren Beruf treiben in wahrem Glauben, mit Gebet und schuldigem Fleiß, 1 Cor. 15, 18.: Eure Arbeit ist nicht vergeblich in dem HErrn. Darauf höret auch an das 28. Kapitel.

\section[Von Ahas, dem Könige in Juda, und seinem]{Von Ahas, dem Könige in Juda, und seinem gottlosen Wesen}
\quellseite{658}

\section{Summarium}

Ist eine Historia, wie die Abgötterei am König Ahas greulich gestraft wird. Erstlich, daß ihm die Syrer ins Land fallen und viele Leute wegführen, V. 5. Darnach der König Israel, Pekah, fällt ihm auch ins Land und erschlägt ihm auf einen Tag hundert und zwanzig tausend Mann, und führet zwei hundert tausend mit sich gefangen weg, V. 6. An solcher Strafe sollte weltliche Obrigkeit und sonst jedermann lernen, wie greuliche Sünde es um Abgötterei, falschen Gottesdienst und unrechte Lehre sei, und derhalben mit höchstem Fleiß sich dafür hüten, und sich zu dem Haufen halten, da rechte Lehre und der Gottesdienst rein und rechtschaffen ist.

Darnach ist eine Predigt des Propheten Oded, in welcher er solche Tyrannei strafet, welche das Königreich Israel an ihren Brüdern geübt hatte, V. 9. Denn Gott will, daß die Strafe nicht gar ohne alle Maße gehen soll.

a Unsere Israel.

b Bloß hieße war, daß das Volk nicht unter Gott nach seinem Worte lebte, sondern frei, nach seinem eigenen Gutdünken im Gottesdienst, wie 2 Mos. 32, 25. Aaron das Volk entblößet.

\section{Summarium}

Dies ist eine Historia, in welcher man siehet, je mehr Gott den König Ahas strafet um seiner
\quellseite{659}
Abgötterei willen, je mehr treibet er Abgötterei. Gleich wie unsere Papisten auch thun: wenn sie Gott strafen mit den Türken und andern Unglück, richten sie Procession an, lassen Messe lesen, suchen Wallfahrten und geloben sich zu den Heiligen. Suchen also wie Ahab, des rechten Gottes Hülfe nicht, sondern setzen alle ihren Trost und Hülfe auf ihre eignen Werke und abgöttischen Ceremonien. Darum gehets auch, wie der Text hie meldet, V. 22., daß die Gottlosen in der Noth des Verzweifelns nur je länger je mehr machen; wie an Ahab zu sehen ist, welcher die Thüren am Hause des HErrn zuschloß, und machet ihm Altäre in allen Winkeln zu Jerusalem, V. 24. Also thun jetzt unsere Papisten auch: das heilige Evangelium lästern und verdammen sie als Ketzerei, wollen derhalben in ihren Landen nicht leiden. Dagegen leiden und dulden sie alle Abgötterei, Sünde und Schande. Sollte aber Gott die nicht strafen und solches rächen?

\section{Votum}

GOtt sei Lob, Ehre und Dank, der uns von Abgötterei erlöset hat, der stehe uns ferner bei; strafe an uns nicht anderer Leute Abgötterei, und helfe, daß ein jedes unter uns, nach seinem Beruf und Vermögen, sein Wort fördern helfe, Amen.

\chapter{Das 29. Kapitel}

\section{Vorrede}

Keinen Gottesdienst erfordert Gott so ernstlich, hart und heftig, als daß man die Abgötterei abschaffen und rechte Lehre und Gottesdienst fördern und erhalten soll. Derwegen alle christliche Obrigkeit die Abgötterei aus den Augen ihrer Unterthanen, die Prediger aber dieselbige aus den Herzen ihrer Zuhörer abschaffen und wegräumen sollen. Hieher gehören diese Sprüche. Von der Obrigkeit sagt das Buch der Weisheit im 6. Kapitel 22.: Wollt ihr nun, ihr Tyrannen, das ist Regenten, in Volk gerne Könige und Fürsten sein; so haltet die Weisheit in Ehren, auf daß ihr ewiglich herrschet. Von den Predigern insgemein lehret auch St. Paulus 1 Tim. 4, 16.: Hab Acht auf dich selbst und auf die Lehre, beharre in diesen Stücken. Denn, wo du solches thust, wirst du dich selbst selig machen, und die dich hören. Solche Obrigkeit und Lehrer hat der treue, barmherzige Gott dieser löblichen Gemeine auch gegeben. Dafür danken wir dir heute alle, o reicher und milder Gott ins Himmel Throne, mit höchstem Fleiß, und bitten dich herzlich, du wollest unsere liebe Obrigkeit und Prediger ferner segnen, stärken und erhalten, und ihres gleichen geben bis an den jüngsten Tag. Dazu sprech alle von Herzen: Amen. Darauf höret ferner an das 29. Kapitel.

\section[Jehiskia richtet den rechten Gottesdienst wieder an]{Jehiskia richtet den rechten Gottesdienst wieder an}
\quellseite{660}

\section{Summarium}

Ist eine sonderliche Historia, welche alle Regenten sollten mit Fleiß bedenken und sie ihnen lassen befohlen sein wie einen Spiegel, darin sie sehen sollten, was fürnehmlich ihr Amt wäre, wenn sie wollten recht und glücklich regieren: nämlich, daß sie am ersten sollten den Tempel öffnen, daß ist, Gottes Wort und den rechten Gottesdienst fördern und unter die Leute bringen, wie der fromme König Hiskiä thut, welcher die Priester und Leviten seine Söhne heißt, sie bittet und vermahnet, sie sollen nicht nachlässig sein, sondern des Gottesdienstes mit Fleiß warten, V. 11. Gegen solchem Könige sollten sich alle gottlose Regenten schämen, so entweder Gottes Wort und desselben Diener verachten, oder wohl zum Land ausjagen, verfolgen und tödten. Aber was für Glück bei ihnen sei, siehet man auch wohl, wer es anders sehen will.
\quellseite{661}

\section{Summarium}

Ist ein schön Exempel, wie der König Hiskiä selbst mit und bei dem Gottesdienst ist. Denn solches ist weltliche Obrigkeit schuldig, auf daß sie mit ihrem Exempel den Unterthanen fürgehe. Daß aber der Herr sagt vom Sündopfer für den König und das ganze Volk, ist nicht zu verstehen, als sei ihnen die Sünde um solches Opfers willen vergeben. Denn Vergebung der Sünden ist je und je gestanden auf dem einigen Opfer des Sohnes Gottes. Desselben Blut allein erwirbt Vergebung der Sünden. Die Sündopfer aber im Gesetz sind nur eine äußerliche Anzeigung gewesen, damit das Volk bekennen sollte, daß sie wider Gottes Gebot und Satzung gehandelt hätten, und also öffentlich solche Mißhandlung für der Gemeine abkehren mit gewöhnlichem Opfer, auf daß sie nicht würden von der Gemeine, als öffentliche Verbrecher, abgesondert. Also lässet der König sich und das Volk versöhnen, zum öffentlichen Bekenntniß, daß sie sich begehreten zu bessern; aber durch solches Opfern ist die Sünde für Gott nicht bezahlet worden, sondern sie haben sich mit dem Herzen zu Gott wenden und Vergebung durch den verheißenen Samen hoffen müssen. Da ist die rechte Vergebung gefolget für Gott, im Gewissen, da jene Vergebung allein für die Leute gegolten hat.

\section{Votum}

GOtt sei uns gnädig und vergebe uns alle unsere Sünde; er wolle auch ferner allzeit geben und erhalten dieser gemeinen Stadt und Vaterland christliche Obrigkeit und treue Lehrer, bis zur Wiederkunft unsers HErrn JEsu Christi, und solches zu seinem Lob und den gehorsamen Unterthanen und Zuhörern zur seligen Wohlfahrt, Amen.

\chapter{Das 30. Kapitel}

\section{Vorrede}

Man schreibt, daß Diagoras Melius, ein Epicurer, der heidnischen Opfer gespottet und gesagt: Es schmecke ihm dennoch Essen und Trinken wohl, ob er gleich den Göttern nicht opferte. Des haben die Athenienser also geeifert, daß er hat entlaufen müssen, und haben sein ganzes Vaterland, die Insel Melos, verwüstet, dieweil ein solch epicurisch Weltkind darinnen erzogen war. Sollte des nun Gott, der HErr, nicht ernstlich eifern und hart strafen, wenn man die seligmachende Lehre seines Worts verachtet? also: ein anderer, mit Namen Pherecides, ist von den Läusen aufgefressen, daß er die heidnischen Opfer verachtete. So wird freilich das höllische Feuer auffressen, die Gottes Wort abermal verachten, darinnen gepredigt wird von dem rechten Opfer für unsere Sünde, nämlich von unsers lieben HErrn JEsu Christi Person, Verdienst, Wohlthaten und Willen. Derwegen wollen wir uns alle hüten für Verachtung göttlichen Wortes, und wollen dagegen dasselbe gern hören, fleißig lernen und christlich darnach glauben und leben, in Kraft und durch Regierung Gottes, des heiligen Geistes, Amen. Folgendes höret an das 30. Kapitel.

\section[Ausrufung und Haltung des Passah]{Ausrufung und Haltung des Passah}
\quellseite{662}

\section{Summarium}

Dies ist eine schöne Historia, wie der fromme König Hiskiä Prediger ausgeschickt in das ganze Land, und vermahnet sein Volk zur Buße, V. 1. Und tröstet sie wiederum, wenn sie sich von Herzen bekehren zum HErrn, so sei Gott gnädig und werde nicht allein ihnen, sondern auch ihren Kindern, so im Gefängniß unter den Heiden waren, gnädige Hülfe erzeigen, V. 6. Solche Predigt und Vermahnung verlacht der meiste Theil, wie denn der Welt Art ist, V. 10. Dennoch waren etliche, die es annahmen und besserten sich, V. 11. Denn Gottes Wort gehet nicht ab ohne Frucht.

Darum soll dies Exempel ernstlich die Obrigkeit wohl bedenken, und wissen, daß sie für Gott schuldig sei, rechte Prediger zu bestellen, und unrechten Gottesdienst abzuthun, wie Hiskiä hie thut.

Zum andern sollen die Unterthanen sich sonderlich des befleißigen, daß sie Gottes Wort nicht verachten, sondern von Herzen und mit großem Ernst es annehmen, und sich bessern.

a Heiligen Reinigkeit) das ist, Gott siehet das Herz an; wenn das rechtschaffen ist an Gott, so fraget er nicht nach äußerlicher Reinigkeit, die nach dem Gesetze heilig ist.
\quellseite{663}

\section{Summarium}

Ist ein Exempel, welches sonderlich wohl zu merken ist. Es hatte seine sonderliche Ordnung mit dem Osterlamm, wie im andern Buch Mose stehet. Nun konnte man in der Eile nicht aller Dinge solche Bräuche und Ceremonien alle halten, wie im Gesetz geboten war.

Da greift der König Hiskia in solche Gesetzordnung, und lehret, man soll mehr darauf sehen, daß das Herz gegen Gott rechtschaffen sei und sich des äußerlichen Wesens halber nicht sonderlich bekümmern.

Denn Gott wolle gnädig sein, nicht um solcher äußerlichen Werke oder Reinigkeit willen, sondern um sein selbst willen, so fern die Herzen rein sind, das ist, so sie ihre Sünde erkennen und auf Gottes Zusagung und Barmherzigkeit trauen; zeigt dadurch an, daß es alles nicht am äußerlichen Thun und Lassen, sondern am Herzen gelegen sei.

\section{Votum}

Gott, der heilige Geist, erleuchte und regiere uns, daß wir innerlich und äußerlich mit Herz, Worten und Thaten fleißig christlich leben und die Gottesdienste üben; und helfe, daß unsere liebe Obrigkeit ferner seine Kirche fördere und schütze, und daß wir Gottes Wort hoch achten, und zu unserer Besserung gebrauchen, Amen.

\chapter{Das 31. Kapitel}

\section{Vorrede}

Gott befehlet der Obrigkeit im 5. Buch Mose im 12. Kapitel V. 3. also: Reißet um die Altäre der Heiden, und die Götzen ihrer Götter thut ab, und vertilget ihren Namen. St. Paulus Gal. 6, 6. lehret dankbare Kirch-Kinder also: Der unterrichtet wird mit dem Wort, der theile mit allerlei Gutes dem, der ihn unterrichtet. Und Sirach im 1. Kapitel V. 19. unterweiset insgemein alle Christen also: Wer den HErrn fürchtet, dem wird wohl gehen. Dieser Sprüche Exempel erzählet das 31. Kapitel.

\section[Abschaffung der Abgötterei, Versorgung der Priester und Leviten]{Abschaffung der Abgötterei, Versorgung der Priester und Leviten}

a Anhalten) auf daß sie nicht aus Mangel der Nahrung müßten die Bücher lassen, beide zu studiren und zu lehren, und ihre Nahrung suchen. Denn Kirchen-Diener sollen versorget sein und studiren, wie Nehemias und Sirach auch sagen.
\quellseite{664}
und ist noch viel überblieben; denn der HErr hat sein Volk gesegnet, darum ist dieser Haufe überblieben. *4 Mos. 5, 9.

\section{Summarium}

Die ist erstlich eine Historia, wie dem Exempel und Befehl des Königs Hiskiä nach alle abgöttischen Altäre und Kirchen im Lande sind umgerissen worden. Denn weltlicher Obrigkeit (wie oft gemeldet) stehets sonderlich zu, daß sie solches nicht dulden, sondern mit Ernst abschaffen sollen.

Darnach ist eine andere Historia, daß der König das Volk vermahnet, sie sollen ihnen die Leviten und Priester, so den Gottesdienst verwalten müssen, in lassen befohlen sein und von dem Ihren mildiglich und treulich ihnen geben, V. 4. Denn weltliche Obrigkeit soll mit allem Fleiß daran sein, daß sie nicht allein selbst gern dazu geben, daß die Kirchen-Diener mögen erhalten werden, sondern auch ihre Unterthanen dahin halten, daß sie es auch thun.

Am Ende ist ein schöner Spruch, nach welchem sich alle Menschen, sonderlich aber weltliche Obrigkeit halten soll, da er spricht: Hiskiä habe in allem, das er anfing, Gott von Herzen gesucht; darum habe er auch Glück gehabt.

\section{Votum}

Regiere und helfe uns Gott, um Jesu Christi willen, durch den heiligen Geist, daß wir alle Abgötterei fliehen, das Predigtamt nach Vermögen helfen fördern, und aus Gnaden Glück, zeitlich und ewig haben, an Leib und Seele, Amen.

\chapter{Das 32. Capitel}

\section{Vorrede}

Der königliche Prophet David sagt im 34. Psalm V. 7.: Da dieser Elende rief, hörte der HErr, und half ihm aus allen seinen Nöthen. Es ein tröstlicher Spruch für alle Nothleidende, die Gott unterm Kreuz in wahrem Glauben anrufen. Aber da möchte jemand gedenken und sagen: Ich erfahre das Gegenspiel, daß manches fleißig und herzlich betet und gleichwohl in seiner Noth stecken bleibt, als in Armuth, Verfolgung, Krankheit und dergleichen. Da sollt ihr wissen, ihr Kinder Gottes, daß wir in zeitlichen Nöthen unsern Willen in unsers lieben Gottes Willen stellen und auf ihn hoffen sollen; ob er uns gleich tödtet, und Leib und Seele ließ verschmachten. Er hat sein Bedenken, warum er also die Seinen mit der Noth ringen läßt; und langsam kömmt er, kömmt aber wohl; und, lange verzogen, heißt nicht gelogen; und im Fall, daß wir hie in diesem Leben (weil Gott weiß, daß es uns nicht zu unserer Seelen Heil und Seligkeit nutz und gut ist) nicht errettet würden; so sollen wir doch reichlich und überflüssig in jenem Leben gewißlich erhöret und unsere Bitte tausend, ja hundert tausendfältig gewähret werden. Denn Gott läßt die Zinsen zusammen über einen Haufen wachsen und wird sie uns darnach auf einmal über einen Haufen geben, im ewigen Leben. Dieses lehret auch das 32. Kapitel.
\quellseite{665}

\section[Hiskias wunderbare Beschützung wider Sanherib]{Hiskias wunderbare Beschützung wider Sanherib}

\section{Summarium}

Ist eine Historia, daß der fromme König Hiskia, der Gottes Wort und Dienst so herrlich gefördert hat, dennoch sammt den Seinen das Kreuz und sehr große Krieges-Noth tragen müssen. Da ist nun sonderlich zu merken, wie er in solcher Noth sein Volk nicht auf Menschen-Hülfe, sondern auf Gottes Hülfe weiset und tröstet, welche Hülfe allein nicht lügen und fehlen kann, V. 7.

Wiederum, was man außer derselbigen Hülfe suchet und hoffet, es sei Geld, Gut, Gewalt, Freundschaft, feste Gebäude, großer Vorrath und anders, ist alles nicht allein ungewisse, sondern auch schädlich, wo man Gott nicht dabei hat, wie alle Historien zeugen.
\quellseite{666}
*Enthielt, wie ein Hirt seine Schafe hält oder die Wölfe, und hütet, daß sie gehen hin und her zur Weide: also konnten die zu Jerusalem auch aus- und einziehen sicher etc.

\section{Summarium}

Dies ist erstlich ein Exempel, daß Gott, wenn man ihn in der Noth sucht und seiner Hülfe gewartet, das Gebet annehmen und gnädig helfen wolle, sollte es gleich durch Wunderzeichen müssen geschehen.

Darnach ist ein ander Exempel, B. 27., in dem wir sehen, wie auch große Heiligen, so durch den festen Glauben in der höchsten Noth bestanden sind, in Glück und Wohlfahrt nicht bestehen und sich durch Sicherheit und Hoffart versündigen, auf daß wir lernen in Gottes Furcht wandeln und täglich bitten, daß Gott die Hand nicht von uns abziehe, sondern uns in der Anfechtung erhalten und behüten wolle.

\section{Votum}

Helfe GOtt die hochgelobte Dreifaltigkeit, daß wir beständig bleiben in wahrem Glauben und in der Gottseligkeit; und erhöre auch unser Gebet, und helfe uns aus aller Noth, Leibes und der Seelen, Amen.

\chapter{Das 33. Kapitel}

\section{Vorrede}

Der heilige Prophet Micha sagt im 7. Kapitel B. 18. 19.: Wo ist solch ein Gott, wie du bist? Der die Sünde vergibt und erläßt die Missethat den Uebrigen seines Erbtheils. Der seinen Zorn nicht ewiglich behält, denn er ist barmherzig. Er wird sich unser wieder erbarmen, unsere Missethat dämpfen und alle unsere Sünde in die Tiefe des Meers werfen. Wenn nun euch Gott, den HErrn, als einen unbarmherzigen Stockmeister der Teufel will einbilden, der sich weder erbarmen, noch Sünde vergeben, sondern nur ewiglich Zorn halten wollte; alsdann sollt ihr ihn abfertigen mit seinem falschen Bilde, und zu ihm sprechen: Gehe hin, Teufel, und lerne das Handwerk daß, du kannst Gott nicht recht abconterfeien, noch seine rechte Gestalt treffen. Es ist ein großer Prophet, der heißt Micha, der hat mir Gott, den Allmächtigen, viel anders fürgemalet und geben ihm alle Propheten Beifall, nämlich, daß Gott Sünde vergebe und erlasse die Missethat allen denen, die sich zu ihm bekehren; er erbarme sich über uns, er dämpfe unsere Missethat, und werfe alle unsere Sünde in die Tiefe des Meers, das ist, er vertilget sie also aus Gnaden, daß er derselben nimmermehr in Ungnaden eingedenk sein will, nicht anders, als wenn er sie alle in ein Gebündlein zusammen geknüpfet und mit einander ins Meer geworfen hätte, daraus sie nimmermehr empor schwimmen, noch ans Licht
\quellseite{667}
kommen sollen. So heißt der Text, Teufel! Darum troll dich, du feindseliger Geist, wohin du gehörst; mir soll Gottes Wort, welches in Ewigkeit bleibet, viel theurer und werther sein, denn alle dein falsches Eingeben. Denn Gott lügt nicht und sein Wort und Verheißung fehlet nicht. Dies betrachtet auch bei dem 33. Kapitel.

\section[Regierung Manasses und Amons]{Regierung Manasses und Amons}
\quellseite{668}

\section{Summarium}

Dies ist eine schöne Historia, in welcher man erstlich siehet, wie der fromme König Hiskiä einen ungerathenen gottlosen Sohn hat, B. 1., der allerlei Abgötterei stiftet und auf Gottes Wort keine Achtung hat. Darum wird er von den Assyrern überfallen und gen Babel gefangen weggeführet, B. 11.

Da ist nun sonderlich ein Exempel der großen Gnade und Barmherzigkeit Gottes, daß er einen solchen gottlosen Tyrannen, nachdem er seine Sünde erkennt und um Gnade bittet, wiederum zu Gnaden annimmt und in die vorigen Würden und Gewalt setzet, B. 13. Zu Trost uns allen, daß wir von Sünden ablassen und uns bessern, und alsdann der Gnade Gottes gewarten sollen und ja nicht verzweifeln.

Sonderlich aber ist zu merken, daß Manasse, nachdem er wieder vom Gefängniß ledig war, alle Abgötterei abschaffet und dem ganzen Volk befiehlet, daß sie dem Herrn, dem Gott Israel dienen sollen, B. 14. Welches eine Anzeigung ist einer rechten Buße \&c.

Am Ende ist ein ander Exempel, von seinem Sohn Amon, der gottlos war, B. 20., aber besserte sich nicht wie sein Vater, darum treibt ers nicht lange. Denn er ward von seinen eigenen Knechten ermordet, B. 24. Aber die Aufrührer wurden gestraft, B. 25.

\section{Gebet}

Der barmherzige Gott vergebe uns auch gnädiglich alle unsere Sünde, lindere uns die wohlverdienten Strafen, helfe, daß wir allzeit christlich leben, und regiere die lieben Kinder, daß sie auch Gott und ihren Eltern gehorsam sein, Amen.

\chapter{Das 34. Capitel}

\section{Vorrede}

Wie der Prophet Esaia, also der Sohn Gottes, unser Herr Jesus Christus, nennt die Tempel oder Kirchen Gotteshäuser oder Bethäuser; darum, daß man die zu Gottes Ehre gebrauchen soll und darin sein Wort lehren und hören, seine Sacramente handeln und gebrauchen, das Gebete und andere Gottesdienste verrichten. Zu solchem Gebrauch der Kirchen- und Gotteshäuser haben allezeit alle christliche Obrigkeit und treue Seelsorger ihre Unterthanen und Zuhörer fleißig ermahnet und gehalten. Denn, anderer Exempel zu geschweigen: Im Concilio Elibertino ist beschlossen worden, wenn einer einheimisch und gesund gewesen, und in drei Wochen in keine Kirche kommen, hat man ihn unwerth geachtet, daß er zum Abendmahl des Herrn gelassen würde. Desgleichen zu Karthago, in vierter Synode, hat man also geschärfet, daß man die in Bann gehen, die unter der Predigt aus der Kirche ohne Noth gelaufen sind. Derwegen habt lieb die Stätte des Hauses Gottes, da man predigt seine Wunder und haltet euch zu seinem Altar, da man höret die Stimme des Dankens, Ps. 26, 7. Und seid fleißig in allen Diensten Gottes. Ferner höret an das 34. Capitel.

\section[Josia reiniget und bessert den Tempel]{Josia reiniget und bessert den Tempel}
\quellseite{669}

\section{Summarium}

Hie ist eine schöne Historia von dem frommen gottseligen König Josia, daß er erstlich alle Abgötterei ausrottet, V. 3., und darnach den Tempel wieder zugerichtet, daß man den Gottesdienst darin üben konnte, V. 8. Solche Exempel sollen weltliche Obrigkeit sonderlich merken, daß sie Gottes Wort und Ehre fördern und aller Abgötterei und falschem Gottesdienst mit Ernst wehren sollen; so will Gott bei ihnen sein, Glück und Heil geben.
\quellseite{670}

\section{Summarium}

Sie ist eine sonderliche Historia, wie Josias, nachdem das Gesetzbuch funden und für ihm gelesen ist worden, seine und des Volks Sünde erkennt, und lässet ihm leid sein, V. 19. Solche Reue und Demuth gefällt Gott wohl. Darum verheißt er, obwohl die Sünden müssen gestraft werden, so wolle er doch die Strafe aufhalten und Josiam in Frieden sterben lassen, V. 27. Solche Verheißung sollte uns auch zur Buße und Besserung bewegen.

Am Ende ist eine Historia, daß Josia das Gesetz-Buch öffentlich habe lesen lassen und sich und das Volk auf ein neues gegen Gott wieder verpflichtet und mit ihm verbunden.

\section{Votum}

GOtt, der Vater unsers HErrn Jesu Christi, verleihe uns auch wahre Demuth und Reue über unsere Sünde, vergebe uns auch dieselbe, und ziehe auf die wohlverdienten Strafen; und helfe, daß wir sein Wort gern hören, und darnach christlich und beständig leben, Amen.

\chapter{Das 35. Kapitel}

\section{Vorrede}

Gott will und begehret es nicht, daß wir ihm etwas umsonst thun sollen. Denn Gott ist nicht ungerecht, daß er vergesse eures Werkes und Arbeit in der Liebe, die ihr beweiset habt in seinem Namen, sagt die Epistel an die Ebräer im 6. Kapitel V. 10. Und der Prophet Maleachi 3, 16. spricht: Der HErr merket und höret, und ist für ihm ein Denkzettel geschrieben für die, so den HErrn fürchten und an seinen Namen gedenken. Und Gott träget seine Gläubige und Gehorsame auf seiner Brust, und drücket sie wie ein Siegel auf sein Herz, im Hohenlied Salomonis 8, 6. Auch lehret also der hocherleuchtete Prophet Esaias und spricht im 3. Kapitel V. 10.: Gott läßt predigen von den Gerechten, daß sie es gut haben, denn sie die Frucht ihrer Werke essen sollen. Sonderlich am jüngsten Tage, sagt St. Paulus Röm. 2, 7., wird Gott einem jeden nach seinen Werken geben, nämlich Preis und Ehre und unvergängliches Wesen denen, die mit Geduld in guten Werken trachten nach dem ewigen Leben. Derwegen wollen wir in Gottesfurcht leben, daß Gott aus Gnaden segne, Glück und Wohlfahrt beschere, schütze und behüte uns alle, und einen jeden in seinem Beruf und Stande, Handel und Wandel, Arbeit und Nahrung, Amen. Darauf höret weiter an das 35. Kapitel.

\section[Josia hält Passah, und kommt um im Streit]{Josia hält Passah, und kommt um im Streit}
% Scan page 671: omitted; Bible text/narrative only, no retained apparatus.
\quellseite{672}

\section{Summarium}

Ist eine Historia, wie der fromme Josias das Osterfest wieder anrichtet und reichlich dazu hilft, B. 1. Denn weltliche Obrigkeit soll sichs etwas kosten lassen, daß der rechte Gottesdienst, und sonderlich das Predigt-Amt, recht bestellet und angerichtet werde.

Am Ende ist eine Historia, wie der fromme König Josias sich in einen unnöthigen Krieg eindränget, da er sonst hätte mögen Ruhe haben, und wird also erschossen, B. 22. Denn es fallen auch Heilige oft. Darum bedarf es desto mehr Betens, daß uns Gott in seinem Schuz erhalten und in der Anfechtung behüten wolle.

\section{Datum}

Helffe GOtt, daß unsere liebe Obrigkeit sich ferner des heiligen Predigtamts treulich annehme, und das fördere und erhalte; und bewahre uns in seinem Gehorsam bis an unser Ende, damit wir aus Fürwitz und muthwillig zu unserm Unglück und Herzeleid nicht Ursache geben, Amen.

\chapter{Das 36. Capitel}

\section{Vorrede}

Der weise Mann Sirach sagt im 10. Capitel B. 8.: Um Gewalt, Unrecht und Geiz willen kommt ein Königreich von einem Volk aufs andere. Dieses wahren Spruchs ist ein öffentlich und klar Exempel das jüdische Königreich. Also, hat der gerechte Gott des israelitischen Volkes, das er zuvor allen Völkern, Land und Leuten, vorgezogen hat, nicht verschonet; darum, daß diese vornehmen Hauptlaster, Abgötterei, Sicherheit, Geiz, Ungerechtigkeit, Gewalt und andere Sünde mehr, wie eine Sündflut eingerissen und mit Gewalt im Schwange gegangen sein, Ezech. 16, 49.: auch so mögen wohl alle für Gottes Zorn und Strafen sich auch hüten, die dergleichen Sünden gethan haben und in solchen Lastern noch leben. Denn Gott läßt sich nicht spotten, spricht St. Paulus an die Galater im 6. Capitel B. 7. Und darum sagt Job im 9. Capitel B. 28.: Ich hab Achtung auf alle meine Werke, denn ich wußte, daß du, HErr, derer nicht schonest, die da sündigen. Dieses lehret auch das 36. und letzte Kapitel.

\section[Von Josias Nachfolgern, und der babylonischen Gefangenschaft]{Von Josias Nachfolgern, und der babylonischen Gefangenschaft}
